\chapter{Complexity — EXP}
%%% generated by the blueprint pipeline from TCSlib.Complexity.ClassNP.EXP
%%%%%%%%%%%%%%%%%%%%%%%%%%%%%%%%%%%%%%%%%%%%%%%%%%%%%%%%%%%%%%%%%%%%%%%%%%%%%%%
\section{Overview}

This module introduces the exponential-time classes: $\mathrm{EXP}$ as the union of
the deterministic time classes $\mathrm{DTIME}(2^{n^c})$, and $\mathrm{NEXP}$ in
certificate form (exponential-length certificates checked by a polynomial-time
verifier language). It proves the chain of inclusions
$\mathrm{P} \subseteq \mathrm{EXP}$, $\mathrm{NP} \subseteq \mathrm{EXP}$ (by
brute-force certificate enumeration on a concrete multi-tape machine), and
$\mathrm{EXP} \subseteq \mathrm{NEXP}$ (by padding), together with the supporting
fixed-width binary counter, reversible logging/restoration, clean-call, and
exponential-split machinery.

%%%%%%%%%%%%%%%%%%%%%%%%%%%%%%%%%%%%%%%%%%%%%%%%%%%%%%%%%%%%%%%%%%%%%%%%%%%%%%%
\section{Declarations}

\begin{definition}[The class $\mathrm{EXP}$]
\lean{Complexity.EXP}
\label{Complexity.EXP}
\leanok
\uses{Complexity.DTIME}
A language $L$ over the binary alphabet belongs to $\mathrm{EXP}$ if and only if
$L \in \mathrm{DTIME}(2^{n^c})$ for some constant $c \in \mathbb{N}$, i.e.\ $L$ is
decided by some finite multi-tape Turing machine within $a \cdot 2^{n^c}$ steps on
inputs of length $n$, for constants $a$ and $c$.
\end{definition}

\begin{definition}[Exponentially bounded function]
\lean{Complexity.ExpBound}
\label{Complexity.ExpBound}
\leanok
A function $p : \mathbb{N} \to \mathbb{N}$ is \emph{exponentially bounded} if there
exist constants $C, c \in \mathbb{N}$ such that $p(n) \le C \cdot 2^{(n+1)^c}$ for
every $n$.
\end{definition}

\begin{definition}[The class $\mathrm{NEXP}$, certificate form]
\lean{Complexity.NEXP}
\label{Complexity.NEXP}
\leanok
\uses{Complexity.P}
A language $L$ over the binary alphabet belongs to $\mathrm{NEXP}$ if there exist
constants $C, c \in \mathbb{N}$ and a verifier language $V \in \mathrm{P}$ such that
for every binary string $x$: $x \in L$ if and only if there is a certificate $u$ of
length exactly $C \cdot 2^{(|x|+1)^c}$ with the concatenation
$x \mathbin{+\!\!+} u$ in $V$.
\end{definition}

\begin{theorem}[$\mathrm{P} \subseteq \mathrm{EXP}$]
\lean{Complexity.P_subset_EXP}
\label{Complexity.P_subset_EXP}
\leanok
\uses{Complexity.DTIME.mono, Complexity.EXP, Complexity.P}
\difficulty{3}
Every language in $\mathrm{P}$ lies in $\mathrm{EXP}$: a language over the binary
alphabet decidable in time polynomial in the input length is also decidable in time
$2^{n^c}$ for some constant $c$, up to a constant factor.
\end{theorem}

\begin{definition}[Little-endian value of a boolean word]
\lean{Complexity.enumValue}
\label{Complexity.enumValue}
\leanok
The value of a boolean word read as a little-endian binary numeral: the empty word
has value $0$, and a word with first bit $b$ and tail $u$ has value
$2 \cdot \mathrm{val}(u) + b$. Leading (high-order) zero bits are allowed.
\end{definition}

\begin{definition}[Fixed-width binary increment]
\lean{Complexity.enumInc}
\label{Complexity.enumInc}
\leanok
The fixed-width increment of a boolean word viewed as a little-endian binary
counter: it returns the incremented word of the same length, or nothing on
overflow (when the word is all ones); the empty word overflows.
\end{definition}

\begin{definition}[Width-$w$ little-endian binary word]
\lean{Complexity.enumWord}
\label{Complexity.enumWord}
\leanok
For all natural numbers $w$ and $i$, the width-$w$ little-endian binary word of
$i$ is the list of the $w$ lowest binary digits of $i$, least significant first,
with high-order zeros retained.
\end{definition}

\begin{lemma}[Value bound for a boolean word]
\lean{Complexity.enumValue_lt}
\label{Complexity.enumValue_lt}
\leanok
\uses{Complexity.enumValue}
\difficulty{2}
For every boolean word $u$, the little-endian value of $u$ is strictly less than
$2^{|u|}$.
\end{lemma}

\begin{lemma}[Length of the width-$w$ word]
\lean{Complexity.enumWord_length}
\label{Complexity.enumWord_length}
\leanok
\uses{Complexity.enumWord}
\difficulty{2}
For all $w$ and $i$, the width-$w$ little-endian binary word of $i$ has length
exactly $w$.
\end{lemma}

\begin{lemma}[The word of rank zero]
\lean{Complexity.enumWord_zero}
\label{Complexity.enumWord_zero}
\leanok
\uses{Complexity.enumWord}
\difficulty{2}
For every width $w$, the width-$w$ little-endian binary word of $0$ is the
all-false word of length $w$.
\end{lemma}

\begin{lemma}[Value of an in-range word]
\lean{Complexity.enumWord_value}
\label{Complexity.enumWord_value}
\leanok
\uses{Complexity.enumValue, Complexity.enumWord}
\difficulty{3}
If $i < 2^w$, then the little-endian value of the width-$w$ binary word of $i$
equals $i$.
\end{lemma}

\begin{lemma}[Equal-length words with equal values coincide]
\lean{Complexity.enumValue_injective}
\label{Complexity.enumValue_injective}
\leanok
\uses{Complexity.enumValue}
\difficulty{4}
Two boolean words of the same length with the same little-endian value are equal
(trailing high-order false bits included).
\end{lemma}

\begin{lemma}[Every word represents its own value]
\lean{Complexity.enumWord_complete}
\label{Complexity.enumWord_complete}
\leanok
\uses{Complexity.enumValue, Complexity.enumValue_injective, Complexity.enumValue_lt, Complexity.enumWord, Complexity.enumWord_length, Complexity.enumWord_value}
\difficulty{1}
Every boolean word $u$ equals the width-$|u|$ little-endian binary word of its own
little-endian value.
\end{lemma}

\begin{lemma}[Specification of the fixed-width increment]
\lean{Complexity.enumInc_spec}
\label{Complexity.enumInc_spec}
\leanok
\uses{Complexity.enumInc, Complexity.enumValue}
\difficulty{4}
For every boolean word $u$: if the fixed-width increment of $u$ succeeds with
result $v$, then $|v| = |u|$ and the little-endian value of $v$ is one more than
that of $u$; if it overflows, then the value of $u$ plus one equals $2^{|u|}$.
\end{lemma}

\begin{lemma}[Increment on canonical words]
\lean{Complexity.enumInc_word}
\label{Complexity.enumInc_word}
\leanok
\uses{Complexity.enumInc, Complexity.enumInc_spec, Complexity.enumValue_injective, Complexity.enumValue_lt, Complexity.enumWord, Complexity.enumWord_length, Complexity.enumWord_value}
\difficulty{4}
For $i < 2^w$, the fixed-width increment of the width-$w$ binary word of $i$ equals
the width-$w$ word of $i+1$ when $i + 1 < 2^w$, and reports overflow otherwise.
This includes width zero.
\end{lemma}

\begin{lemma}[Exact-width certificates versus in-range ranks]
\lean{Complexity.enumCandidates_iff}
\label{Complexity.enumCandidates_iff}
\leanok
\uses{Complexity.enumValue, Complexity.enumValue_lt, Complexity.enumWord, Complexity.enumWord_complete, Complexity.enumWord_length}
\difficulty{3}
For every width $w$ and every predicate $p$ on boolean words: there exists a word
$u$ with $|u| = w$ satisfying $p$ if and only if there exists a rank $i < 2^w$ such
that $p$ holds of the width-$w$ binary word of $i$.
\end{lemma}

\begin{definition}[Boolean test of consecutive ranks]
\lean{Complexity.enumAny}
\label{Complexity.enumAny}
\leanok
Given an acceptance function on natural-number ranks, a start rank $i$, and a
count: the Boolean that is true exactly when one of the $\mathrm{count}$
consecutive ranks starting at $i$ accepts, evaluated by testing the ranks in
order.
\end{definition}

\begin{lemma}[Characterisation of the rank test]
\lean{Complexity.enumAny_iff}
\label{Complexity.enumAny_iff}
\leanok
\uses{Complexity.enumAny}
\difficulty{4}
Testing $\mathrm{count}$ consecutive ranks starting at $i$ against an acceptance
function returns true if and only if some rank $j$ with
$i \le j < i + \mathrm{count}$ accepts.
\end{lemma}

\begin{lemma}[Rank test equals certificate existence]
\lean{Complexity.enumAny_certificates}
\label{Complexity.enumAny_certificates}
\leanok
\uses{Complexity.enumAny, Complexity.enumAny_iff, Complexity.enumCandidates_iff, Complexity.enumWord, Turing.MultiTapeTM.indicator}
\difficulty{2}
Let $V$ be a language over the binary alphabet, $x$ a string, and $w$ a width.
Testing all $2^w$ ranks with the acceptance function ``the indicator of $V$ accepts
$x$ concatenated with the width-$w$ binary word of the rank'' returns true if and
only if there exists $u$ with $|u| = w$ and $x \mathbin{+\!\!+} u \in V$.
\end{lemma}

\begin{lemma}[Timed loop invariant for candidate rounds]
\lean{Complexity.enumLoop_run}
\label{Complexity.enumLoop_run}
\leanok
\uses{Complexity.enumAny, Turing.Cfg, Turing.FinTM, Turing.MultiTapeTM.runFrom, Turing.MultiTapeTM.runFrom_add}
\difficulty{5}
Let $M$ be a finite multi-tape Turing machine over the binary alphabet, $x$ an
input, and let $\mathrm{cfg}$ assign to each rank a configuration of $M$ on $x$.
Suppose the configuration at rank $i + \mathrm{count}$ is halted with single-bit
output false, and that for each rank $j$ with $i \le j < i + \mathrm{count}$ the
machine, within $B$ steps from $\mathrm{cfg}(j)$, either halts with single-bit
output true (when rank $j$ accepts) or reaches $\mathrm{cfg}(j+1)$ (when it
rejects). Then within $\mathrm{count} \cdot B$ steps from $\mathrm{cfg}(i)$ the
machine halts with single-bit output equal to the disjunction of the acceptance
bits over the ranks $i \le j < i + \mathrm{count}$.
\end{lemma}

\begin{lemma}[Polynomial exponents absorb into $2^{n^e}$]
\lean{Complexity.enumExponent_bound}
\label{Complexity.enumExponent_bound}
\leanok
\difficulty{5}
For all constants $K, k \in \mathbb{N}$ there exist constants $A, e \in \mathbb{N}$
such that $2^{K (n+1)^k} \le A \cdot 2^{n^e}$ for every $n$.
\end{lemma}

\begin{lemma}[Enumeration budgets are exponential budgets]
\lean{Complexity.enumBudget_bound}
\label{Complexity.enumBudget_bound}
\leanok
\uses{Complexity.enumExponent_bound}
\difficulty{5}
For all constants $a, C, c, d \in \mathbb{N}$ there exist $A, e \in \mathbb{N}$
such that for every $n$,
\[
a \cdot 2^{C(n+1)^c} \cdot \bigl(n + C(n+1)^c + 1\bigr)^d \le A \cdot 2^{n^e}.
\]
\end{lemma}

\begin{lemma}[Agreement of the two fixed-width increments]
\lean{Complexity.enumCont_inc_eq}
\label{Complexity.enumCont_inc_eq}
\leanok
\uses{Complexity.enumInc, Turing.incFixed}
\difficulty{2}
For every boolean word $s$, the fixed-width little-endian binary increment of $s$
computed by the Turing-machine catalog's counter equals the one computed by the
certificate-enumeration counter: both return the same-length incremented word, or
both fail exactly on overflow (including on the empty word).
\end{lemma}

\begin{lemma}[The stalled increment preserves width]
\lean{Complexity.enumCont_step_length}
\label{Complexity.enumCont_step_length}
\leanok
\uses{Complexity.enumCont_inc_eq, Complexity.enumInc, Complexity.enumInc_spec, Turing.incFixed}
\difficulty{2}
For every boolean word $s$, applying the fixed-width increment and falling back to
$s$ itself on overflow yields a word whose length is exactly $|s|$.
\end{lemma}

\begin{lemma}[Orbit of the stalled increment]
\lean{Complexity.enumCont_orbit}
\label{Complexity.enumCont_orbit}
\leanok
\uses{Complexity.enumCont_inc_eq, Complexity.enumInc_word, Complexity.enumWord, Complexity.enumWord_zero, Turing.incFixed}
\difficulty{3}
For every width $w$ and every $i < 2^w$, iterating the stalled fixed-width
increment (increment, keeping the word unchanged on overflow) $i$ times from the
all-false word of length $w$ yields the width-$w$ little-endian binary word of
$i$. Nothing is asserted at the terminal rank $2^w$.
\end{lemma}

\begin{lemma}[Binary digits of $2^w - 1$]
\lean{Complexity.enumCont_fuel_bits}
\label{Complexity.enumCont_fuel_bits}
\leanok
\difficulty{3}
For every $w$, the binary-digit list of $2^w - 1$ is the all-true word of length
$w$.
\end{lemma}

\begin{definition}[Sparse tape of a finite cell list]
\lean{Complexity.enumCont_sparse}
\label{Complexity.enumCont_sparse}
\leanok
The tape function determined by a finite list of possibly blank cells: a
nonnegative position $z$ holds the $z$-th list entry (which may itself be blank),
and every position beyond the list, as well as every negative position, is blank.
\end{definition}

\begin{lemma}[Appending one sparse cell]
\lean{Complexity.enumCont_sparse_append}
\label{Complexity.enumCont_sparse_append}
\leanok
\uses{Complexity.enumCont_sparse}
\difficulty{4}
Appending one possibly blank entry to a finite cell list changes the induced
sparse tape at exactly the position just past the old list, where the new entry
is written.
\end{lemma}

\begin{lemma}[Erasing the last sparse cell]
\lean{Complexity.enumCont_sparse_erase}
\label{Complexity.enumCont_sparse_erase}
\leanok
\uses{Complexity.enumCont_sparse, Complexity.enumCont_sparse_append}
\difficulty{3}
Let $w$ be a finite list of possibly blank cells and $b$ a possibly blank entry.
Blanking position $|w|$ of the sparse tape of $w$ extended by $b$ yields exactly
the sparse tape of $w$; this holds even when $b$ is itself blank.
\end{lemma}

\begin{definition}[Encoding of head moves as tape symbols]
\lean{Complexity.enumCont_moveCode}
\label{Complexity.enumCont_moveCode}
\leanok
The move code is the encoding of the three head moves as tape symbols: it sends a
left move to a false bit, staying put to the blank symbol, and a right move to a
true bit.
\end{definition}

\begin{definition}[Decoding of the inverse head move]
\lean{Complexity.enumCont_unmove}
\label{Complexity.enumCont_unmove}
\leanok
The inverse-move decoding is the map from recorded move symbols to head moves,
used by the backward restoration pass, that sends a false bit to a right move, the
blank symbol to staying put, and a true bit to a left move; thus every move symbol
is decoded to the inverse of the move it encodes.
\end{definition}

\begin{lemma}[A move and its decoded inverse cancel]
\lean{Complexity.enumCont_unmove_cast}
\label{Complexity.enumCont_unmove_cast}
\leanok
\uses{Complexity.enumCont_moveCode, Complexity.enumCont_unmove}
\difficulty{1}
For every head move $d$, encoding $d$ as a tape symbol and decoding the inverse
move yields exactly the negation of $d$'s integer head displacement.
\end{lemma}

\begin{definition}[History entry of one machine step]
\lean{Complexity.EnumContEntry}
\label{Complexity.EnumContEntry}
\leanok
For a machine with $k$ work tapes, a history entry for one step is a pair
consisting of the tuple of the $k$ symbols under the work heads before the step
and the tuple of the $k$ head moves performed by the step.
\end{definition}

\begin{definition}[Logging instrumentation of one action]
\lean{Complexity.enumCont_logAction}
\label{Complexity.enumCont_logAction}
\leanok
\uses{Complexity.enumCont_moveCode, Turing.Action, Turing.FinTM.tapeBlocks}
Given the symbols currently under the $k$ work heads and an action of a $k$-tape
machine, the corresponding action of a machine with $k + (1 + (k + k))$ work
tapes: it performs the original action on the first $k$ tapes while advancing a
unary clock tape and recording, on $2k$ history tapes, the overwritten symbols and
the encoded head moves. Output and control are otherwise unchanged.
\end{definition}

\begin{definition}[The logging machine]
\lean{Complexity.enumCont_logTM}
\label{Complexity.enumCont_logTM}
\leanok
\uses{Complexity.enumCont_logAction, Turing.FinTM}
For a finite machine $M$ with $k$ work tapes, the logging machine of $M$ is the
machine with the same finite state set and initial state and with
$k + (1 + (k + k))$ work tapes, every transition of which performs $M$'s
transition on the first $k$ tapes while logging a clock mark, the overwritten
symbols, and the head moves on the extra tapes.
\end{definition}

\begin{definition}[Logged configuration]
\lean{Complexity.enumCont_logCfg}
\label{Complexity.enumCont_logCfg}
\leanok
\uses{Complexity.EnumContEntry, Complexity.enumCont_moveCode, Complexity.enumCont_sparse, Turing.Cfg, Turing.FinTM.bufferTape, Turing.FinTM.tapeBlocks}
The configuration of the logging machine corresponding to a source configuration
$c$ and a finite history list $h$: the source fields occupy the first block of
tapes, the clock tape holds $|h|$ true marks, the history tapes hold the recorded
old symbols and encoded moves as sparse tapes, and every history head sits one
cell past the recorded entries.
\end{definition}

\begin{lemma}[One logged action appends one history entry]
\lean{Complexity.enumCont_log_apply}
\label{Complexity.enumCont_log_apply}
\leanok
\uses{Complexity.EnumContEntry, Complexity.enumCont_logAction, Complexity.enumCont_logCfg, Complexity.enumCont_sparse_append, Turing.Action, Turing.Action.apply, Turing.Cfg, Turing.Cfg.workTapeSymbols, Turing.FinTM.bufferTape_append, Turing.FinTM.tapeBlocks}
\difficulty{4}
For any configuration $c$ with history $h$ and any action $a$ of the source
machine: applying the logged version of $a$ to the logged configuration of $c$
with history $h$ yields the logged configuration of $a$ applied to $c$, with
history $h$ extended by the single entry recording the overwritten symbols and
head moves of $a$. This includes halting and output-emitting actions.
\end{lemma}

\begin{definition}[History of a run]
\lean{Complexity.enumCont_history}
\label{Complexity.enumCont_history}
\leanok
\uses{Complexity.EnumContEntry, Turing.Cfg, Turing.Cfg.inputSymbol, Turing.Cfg.workTapeSymbols, Turing.MultiTapeTM, Turing.MultiTapeTM.runFrom}
For a machine, a start configuration, and a time $t$: the list of history entries
(overwritten symbols and head moves) of the actions actually executed during the
first $t$ steps of the run. No entry is recorded once the machine has halted.
\end{definition}

\begin{lemma}[Logging runs in lockstep with the source]
\lean{Complexity.enumCont_log_run}
\label{Complexity.enumCont_log_run}
\leanok
\uses{Complexity.enumCont_history, Complexity.enumCont_logAction, Complexity.enumCont_logCfg, Complexity.enumCont_logTM, Complexity.enumCont_log_apply, Complexity.enumCont_undoTM, Turing.Action.apply, Turing.Cfg, Turing.Cfg.workTapeSymbols, Turing.FinTM, Turing.FinTM.tapeBlocks, Turing.MultiTapeTM.runFrom, Turing.MultiTapeTM.runFrom_succ_eq_step', Turing.MultiTapeTM.step}
\difficulty{5}
For every finite machine $M$, start configuration $c_0$ (possibly prepared, not
necessarily initial), and time $t$: running the logging machine of $M$ for $t$
steps from the logged configuration of $c_0$ with empty history yields the logged
configuration of $M$'s time-$t$ configuration together with the history of the
first $t$ steps. In particular the source's input, output, and halting time are
unchanged by the instrumentation.
\end{lemma}

\begin{definition}[The restoration machine]
\lean{Complexity.enumCont_undoTM}
\label{Complexity.enumCont_undoTM}
\leanok
\uses{Complexity.enumCont_unmove, Turing.FinTM, Turing.FinTM.tapeBlocks}
The restoration controller on $k + (1 + (k + k))$ work tapes: reading the newest
history entry, it alternates moving the $k$ source heads by the inverses of the
recorded moves with writing back the recorded old symbols, erasing the history as
it goes, and halts with all history heads at zero when the left blank of the
clock tape is reached. The native input head never moves.
\end{definition}

\begin{definition}[Restoration checkpoint configuration]
\lean{Complexity.enumCont_undoCfg}
\label{Complexity.enumCont_undoCfg}
\leanok
\uses{Complexity.EnumContEntry, Complexity.enumCont_logCfg, Complexity.enumCont_undoTM, Turing.Cfg, Turing.FinTM.tapeBlocks}
The checkpoint configuration of the restoration machine determined by a source
configuration $c$, a history list $h$, and an input-head position: the source
work fields are those of $c$, the logged history tapes hold $h$, and every
history head inspects the last remaining entry.
\end{definition}

\begin{definition}[Completed restoration configuration]
\lean{Complexity.enumCont_undoResult}
\label{Complexity.enumCont_undoResult}
\leanok
\uses{Complexity.enumCont_logCfg, Complexity.enumCont_undoTM, Turing.Cfg, Turing.FinTM.tapeBlocks}
For a source configuration $c$ and an input-head position $p$, the completed
restoration configuration is the halted configuration of the restoration machine
whose input head is at $p$, whose source work tapes and head positions are those
of $c$, whose clock and history tapes are all blank with heads at zero, and whose
output is empty.
\end{definition}

\begin{lemma}[Undoing one tape write]
\lean{Complexity.enumCont_restore_cell}
\label{Complexity.enumCont_restore_cell}
\leanok
\uses{Turing.Action, Turing.Action.apply, Turing.Action.apply_workTapes, Turing.Cfg, Turing.Cfg.workTapeSymbols}
\difficulty{2}
For any configuration $c$, action $a$, and work-tape index $i$: writing the symbol
that was under head $i$ in $c$, at $c$'s head position, into tape $i$ after $a$
has been applied restores tape $i$ exactly to its contents in $c$. This covers
no-write actions and writes of the blank symbol.
\end{lemma}

\begin{lemma}[Erasing the newest clock mark]
\lean{Complexity.enumCont_clock_erase}
\label{Complexity.enumCont_clock_erase}
\leanok
\uses{Turing.FinTM.bufferTape, Turing.FinTM.bufferTape_append}
\difficulty{3}
For every $n$, blanking cell $n$ of the buffer tape holding $n + 1$ consecutive
true marks yields the buffer tape holding $n$ consecutive true marks.
\end{lemma}

\begin{definition}[Mid-restoration configuration]
\lean{Complexity.enumCont_undoMid}
\label{Complexity.enumCont_undoMid}
\leanok
\uses{Complexity.EnumContEntry, Complexity.enumCont_logCfg, Complexity.enumCont_undoTM, Turing.Action, Turing.Action.apply, Turing.Cfg, Turing.Cfg.workTapeSymbols, Turing.FinTM.tapeBlocks}
The intermediate restoration configuration between inverse movement and inverse
writing, for a source configuration $c$, an action $a$, and a history $h$: the
old symbols of $c$ are retained in finite control, the source heads are back at
their pre-action positions, and the last history entry is already erased.
\end{definition}

\begin{lemma}[First restoration transition: reverse the moves]
\lean{Complexity.enumCont_undo_back}
\label{Complexity.enumCont_undo_back}
\leanok
\uses{Complexity.EnumContEntry, Complexity.enumCont_clock_erase, Complexity.enumCont_logCfg, Complexity.enumCont_moveCode, Complexity.enumCont_sparse, Complexity.enumCont_sparse_erase, Complexity.enumCont_undoCfg, Complexity.enumCont_undoMid, Complexity.enumCont_undoTM, Complexity.enumCont_unmove_cast, Turing.Action, Turing.Action.apply, Turing.Cfg, Turing.Cfg.workTapeSymbols, Turing.FinTM.tapeBlocks, Turing.MultiTapeTM.step, Turing.moveInputPos_zero}
\difficulty{5}
Let $c$ be a source configuration, $a$ an action, $h$ a history list, and $p$ an
input-head position. One step of the restoration machine from the checkpoint at
$p$ for the configuration obtained by applying $a$ to $c$, whose history is $h$
extended by the entry recording $a$ (the symbols under the work heads of $c$ and
the head moves of $a$), equals the mid-restoration configuration for $c$, $a$, $h$
and $p$: the recorded head moves are reversed, the overwritten symbols are
retained in finite control, and the cells of that history entry are erased.
\end{lemma}

\begin{lemma}[Second restoration transition: rewrite the old symbols]
\lean{Complexity.enumCont_undo_write}
\label{Complexity.enumCont_undo_write}
\leanok
\uses{Complexity.EnumContEntry, Complexity.enumCont_logCfg, Complexity.enumCont_restore_cell, Complexity.enumCont_undoCfg, Complexity.enumCont_undoMid, Complexity.enumCont_undoTM, Turing.Action, Turing.Action.apply, Turing.Cfg, Turing.Cfg.workTapeSymbols, Turing.FinTM.tapeBlocks, Turing.MultiTapeTM.step, Turing.moveInputPos_zero}
\difficulty{4}
Let $c$ be a source configuration, $a$ an action, $h$ a history list, and $p$ an
input-head position. One step of the restoration machine from the mid-restoration
configuration for $c$, $a$, $h$ and $p$ equals the checkpoint configuration for $c$
and $h$ at $p$: the retained old symbols are written back, restoring the source
tapes to their contents in $c$, and the history heads move back onto the last
entry of $h$.
\end{lemma}

\begin{lemma}[Final restoration transition on empty history]
\lean{Complexity.enumCont_undo_empty}
\label{Complexity.enumCont_undo_empty}
\leanok
\uses{Complexity.enumCont_logCfg, Complexity.enumCont_undoCfg, Complexity.enumCont_undoResult, Complexity.enumCont_undoTM, Turing.Action.apply, Turing.Cfg, Turing.Cfg.workTapeSymbols, Turing.FinTM.tapeBlocks, Turing.MultiTapeTM.step, Turing.moveInputPos_zero}
\difficulty{4}
With no history left, one silent step of the restoration machine from the
checkpoint for a source configuration $c$ with empty history returns the history
heads to zero and halts, reaching the completed restoration configuration for
$c$. This also covers a source run of zero steps.
\end{lemma}

\begin{lemma}[Complete reversal of a live source prefix]
\lean{Complexity.enumCont_undo_run}
\label{Complexity.enumCont_undo_run}
\leanok
\uses{Complexity.enumCont_history, Complexity.enumCont_undoCfg, Complexity.enumCont_undoResult, Complexity.enumCont_undoTM, Complexity.enumCont_undo_back, Complexity.enumCont_undo_empty, Complexity.enumCont_undo_write, Turing.Action.apply, Turing.Cfg, Turing.Cfg.Halted, Turing.Cfg.inputSymbol, Turing.Cfg.workTapeSymbols, Turing.FinTM, Turing.MultiTapeTM.runFrom, Turing.MultiTapeTM.runFrom_add, Turing.MultiTapeTM.runFrom_succ_eq_step', Turing.MultiTapeTM.runFrom_zero, Turing.MultiTapeTM.step}
\difficulty{5}
Let $M$ be a finite machine that is live (not halted) at every time before $t$
when run from $c_0$. Then the restoration machine, started at the checkpoint for
$M$'s time-$t$ configuration with the $t$-step history, reaches the completed
restoration configuration of $c_0$ --- original work tapes and head positions,
blank history --- in exactly $2t + 1$ steps, preserving any supplied input-head
position.
\end{lemma}

\begin{lemma}[First halting time]
\lean{Complexity.enumCont_first_halt}
\label{Complexity.enumCont_first_halt}
\leanok
\uses{Turing.Cfg, Turing.Cfg.Halted, Turing.MultiTapeTM, Turing.MultiTapeTM.runFrom, Turing.MultiTapeTM.runFrom_add, Turing.MultiTapeTM.runFrom_of_halt}
\difficulty{4}
If a machine's run from $c_0$ is halted at time $T$, then there is a time
$t \le T$ at which it halts for the first time: the machine is live at every time
before $t$, and its configuration at time $t$ equals its configuration at time
$T$.
\end{lemma}

\begin{definition}[Two-head administrative action]
\lean{Complexity.enumCont_bufferAction}
\label{Complexity.enumCont_bufferAction}
\leanok
\uses{Turing.Action}
Given displacements $m$ and $d$ and a next control state, the two-head
administrative action is the action on $L + 1$ work tapes that writes nothing,
emits nothing, moves the native input head by $m$ and the head of the last
(capture) tape by $d$, leaves the heads of the first $L$ work tapes in place, and
selects the given next control state.
\end{definition}

\begin{definition}[Entry action into restoration]
\lean{Complexity.enumCont_undoEntry}
\label{Complexity.enumCont_undoEntry}
\leanok
\uses{Complexity.enumCont_undoTM, Turing.Action, Turing.FinTM.tapeBlocks}
The restoration entry action is the silent action that moves every history head
of the restoration machine, including the clock head, one cell left, from the right
blank onto the newest recorded entry; it writes nothing, leaves the source heads
and the input head fixed, and selects the restoration machine's checkpoint state.
\end{definition}

\begin{definition}[Clean-call wrapper of a machine]
\lean{Complexity.enumCont_cleanTM}
\label{Complexity.enumCont_cleanTM}
\leanok
\uses{Complexity.enumCont_bufferAction, Complexity.enumCont_logTM, Complexity.enumCont_undoEntry, Complexity.enumCont_undoTM, Turing.FinTM, Turing.FinTM.controlAction, Turing.captureAction}
The clean subroutine built from a finite machine $M$: it simulates the logging
version of $M$ with the output captured on one extra final tape, then runs the
captured restoration of all source work, rewinds the native input head and the
captured word, and halts silently. The retained output word lives on the last
work tape; no physical output is emitted.
\end{definition}

\begin{definition}[Administrative configuration of the clean wrapper]
\lean{Complexity.enumCont_cleanCfg}
\label{Complexity.enumCont_cleanCfg}
\leanok
\uses{Complexity.enumCont_cleanTM, Complexity.enumCont_logCfg, Complexity.enumCont_logTM, Turing.Cfg, Turing.FinTM, Turing.FinTM.bufferTape}
An administrative configuration of the clean-call wrapper: the source block holds
the restored work fields of a given source configuration with blank histories,
the last tape holds a retained word $y$, and the control state, native input-head
position, and capture-head position are parameters. The output is empty.
\end{definition}

\begin{lemma}[Captured simulation phase of the clean wrapper]
\lean{Complexity.enumCont_clean_capture}
\label{Complexity.enumCont_clean_capture}
\leanok
\uses{Complexity.enumCont_cleanTM, Complexity.enumCont_concatTM, Complexity.enumCont_history, Complexity.enumCont_logCfg, Complexity.enumCont_logTM, Complexity.enumCont_log_run, Turing.Cfg, Turing.Cfg.Halted, Turing.FinTM, Turing.MultiTapeTM.runFrom, Turing.captureCfg, Turing.capture_run}
\difficulty{3}
If a source machine $M$ is live for $t$ steps from a configuration $c_0$, then the
clean wrapper of $M$, run for $t$ steps from the captured logged start of $c_0$,
reaches the captured logged configuration of $M$'s time-$t$ configuration with the
$t$-step history; any halting emission of $M$ is retained on the capture tape.
\end{lemma}

\begin{lemma}[Dispatch into the captured restoration]
\lean{Complexity.enumCont_clean_undo_entry}
\label{Complexity.enumCont_clean_undo_entry}
\leanok
\uses{Complexity.EnumContEntry, Complexity.enumCont_cleanTM, Complexity.enumCont_concatTM, Complexity.enumCont_logCfg, Complexity.enumCont_logTM, Complexity.enumCont_undoCfg, Complexity.enumCont_undoEntry, Complexity.enumCont_undoTM, Turing.Action.apply, Turing.Cfg, Turing.FinTM, Turing.FinTM.tapeBlocks, Turing.MultiTapeTM.step, Turing.captureAction, Turing.captureCfg, Turing.moveInputPos_zero}
\difficulty{5}
If the source configuration $c$ is halted, then one silent step of the clean
wrapper from the captured logged configuration of $c$ with history $h$ parks
every history head on the last entry and enters the captured restoration phase at
the checkpoint for $c$ and $h$, retaining $c$'s output on the capture tape.
\end{lemma}

\begin{lemma}[Captured restoration returns within $2t + 1$ steps]
\lean{Complexity.enumCont_clean_restore}
\label{Complexity.enumCont_clean_restore}
\leanok
\uses{Complexity.enumCont_cleanCfg, Complexity.enumCont_cleanTM, Complexity.enumCont_concatTM, Complexity.enumCont_first_halt, Complexity.enumCont_history, Complexity.enumCont_logCfg, Complexity.enumCont_logTM, Complexity.enumCont_undoCfg, Complexity.enumCont_undoResult, Complexity.enumCont_undoTM, Complexity.enumCont_undo_run, Turing.Cfg, Turing.Cfg.Halted, Turing.FinTM, Turing.MultiTapeTM.runFrom, Turing.captureCfg, Turing.capture_run}
\difficulty{4}
If the source $M$ is live for $t$ steps from $c_0$, then within $2t + 1$ steps the
clean wrapper's captured restoration, started at the checkpoint of the time-$t$
configuration with its history, reaches the administrative configuration in which
the original source work fields of $c_0$ are restored, the time-$t$ output is
retained on the last tape, and the capture head stands at the end of that word.
\end{lemma}

\begin{lemma}[Effect of a two-head administrative action]
\lean{Complexity.enumCont_buffer_apply}
\label{Complexity.enumCont_buffer_apply}
\leanok
\uses{Complexity.enumCont_bufferAction, Complexity.enumCont_cleanCfg, Complexity.enumCont_cleanTM, Complexity.enumCont_concatTM, Complexity.enumCont_logTM, Turing.Action.apply, Turing.Cfg, Turing.FinTM, Turing.moveInputPos}
\difficulty{3}
For any displacements $m, d$ and control states $q, q'$, applying the two-head
administrative action with displacements $m, d$ and next state $q'$ to an
administrative configuration of the clean wrapper in state $q$ yields the
administrative configuration with the same tapes, retained word and empty output,
in state $q'$, with the native input head moved by $m$ and the capture head moved
by $d$.
\end{lemma}

\begin{lemma}[Bounded input rewind]
\lean{Complexity.enumCont_rewind}
\label{Complexity.enumCont_rewind}
\leanok
\uses{Complexity.enumCont_concatTM, Turing.Cfg, Turing.FinTM.controlAction, Turing.FinTM.controlAction_apply, Turing.FinTM.moveInputPos_neg_val, Turing.FinTM.rewind_scan, Turing.MultiTapeTM, Turing.MultiTapeTM.runFrom, Turing.MultiTapeTM.runFrom_add, Turing.MultiTapeTM.step, Turing.moveInputPos}
\difficulty{4}
Consider a machine with a start state whose transition unconditionally steps the
input head left into a scanning state, and a scanning state that keeps stepping
left over input symbols and exits right into a destination state at the left
blank. Then from any configuration in the start state there is a number of steps,
at most the current input position plus $2$, after which the configuration equals
the original one with control transferred to the destination state and the input
head at cell one.
\end{lemma}

\begin{lemma}[Rewinding the retained output word]
\lean{Complexity.enumCont_clean_buffer_rewind}
\label{Complexity.enumCont_clean_buffer_rewind}
\leanok
\uses{Complexity.enumCont_buffer_apply, Complexity.enumCont_cleanCfg, Complexity.enumCont_cleanTM, Complexity.enumCont_concatTM, Complexity.enumCont_logTM, Turing.Action.apply, Turing.Cfg, Turing.Cfg.workTapeSymbols, Turing.FinTM, Turing.MultiTapeTM.runFrom, Turing.MultiTapeTM.runFrom_succ_eq_step, Turing.MultiTapeTM.runFrom_succ_eq_step', Turing.MultiTapeTM.runFrom_zero, Turing.MultiTapeTM.step}
\difficulty{5}
For every $j \le |y|$: from the administrative configuration of the clean wrapper
with retained word $y$ and capture head just left of cell $j$, exactly $j + 1$
steps return the capture head to zero without erasing $y$, advancing control to
the final administrative state.
\end{lemma}

\begin{lemma}[A halting call can be made clean]
\lean{Complexity.enumCont_clean_complete}
\label{Complexity.enumCont_clean_complete}
\leanok
\uses{Complexity.enumCont_bufferAction, Complexity.enumCont_buffer_apply, Complexity.enumCont_cleanCfg, Complexity.enumCont_cleanTM, Complexity.enumCont_clean_buffer_rewind, Complexity.enumCont_clean_capture, Complexity.enumCont_clean_restore, Complexity.enumCont_clean_undo_entry, Complexity.enumCont_concatTM, Complexity.enumCont_first_halt, Complexity.enumCont_history, Complexity.enumCont_logCfg, Complexity.enumCont_rewind, Turing.Action.apply, Turing.Cfg, Turing.FinTM, Turing.FinTM.controlAction, Turing.FinTM.controlAction_apply, Turing.MultiTapeTM.runFrom, Turing.MultiTapeTM.runFrom_add, Turing.MultiTapeTM.runFrom_succ_eq_step', Turing.MultiTapeTM.step, Turing.captureCfg, Turing.moveInputPos_zero}
\difficulty{6}
If a machine $M$ halts within $T$ steps from configuration $c_0$ on input $x$ with
output $y$, then within $3T + |x| + |y| + 8$ steps the clean wrapper of $M$,
started at the captured logged start of $c_0$, halts in the administrative
configuration with $c_0$'s work fields restored, all histories blank, $y$
retained on the last tape, the input head at cell one, the capture head at zero,
and no physical output.
\end{lemma}

\begin{definition}[The assembly machine]
\lean{Complexity.enumCont_concatTM}
\label{Complexity.enumCont_concatTM}
\leanok
\uses{Turing.FinTM, Turing.FinTM.controlAction}
A two-state machine with one work tape that first emits (copies to its output)
the entire native input and then emits the word written on its work tape, never
writing that tape, and halts at the word's right blank.
\end{definition}

\begin{definition}[Assembly configurations]
\lean{Complexity.enumCont_concatCfg}
\label{Complexity.enumCont_concatCfg}
\leanok
\uses{Turing.Cfg, Turing.FinTM.bufferTape}
An assembly configuration is the configuration of the assembly machine on input
$x$ determined by a control state, an input-head position, a candidate word $s$
held fixed on the single work tape with its head at a given cell, and the output
prefix emitted so far.
\end{definition}

\begin{lemma}[Input phase of assembly]
\lean{Complexity.enumCont_concat_native}
\label{Complexity.enumCont_concat_native}
\leanok
\uses{Complexity.enumCont_concatCfg, Complexity.enumCont_concatTM, Turing.Action.apply, Turing.Cfg.inputSymbol, Turing.FinTM.controlAction_apply, Turing.FinTM.inputSymbol_at, Turing.MultiTapeTM.runFrom, Turing.MultiTapeTM.runFrom_succ_eq_step, Turing.MultiTapeTM.runFrom_succ_eq_step', Turing.MultiTapeTM.runFrom_zero, Turing.MultiTapeTM.step, Turing.moveInputPos_pos_of_ne_right, Turing.moveInputPos_zero}
\difficulty{5}
Suppose the input $x$ decomposes as $\mathit{pre} \mathbin{+\!\!+} \mathit{rest}$,
and let $s$ be a candidate word. After exactly $|\mathit{rest}| + 1$ steps from the
input-phase configuration of the assembly machine with the input head on the first
symbol of $\mathit{rest}$, candidate $s$ with its head at zero, and output
$\mathit{pre}$, the machine is in the candidate phase with the input head at the
right boundary, the candidate tape and its head at zero unchanged, and output $x$.
\end{lemma}

\begin{lemma}[Candidate phase of assembly]
\lean{Complexity.enumCont_concat_candidate}
\label{Complexity.enumCont_concat_candidate}
\leanok
\uses{Complexity.enumCont_concatCfg, Complexity.enumCont_concatTM, Turing.Action.apply, Turing.Cfg.workTapeSymbols, Turing.FinTM.bufferTape_nat, Turing.FinTM.controlAction_apply, Turing.MultiTapeTM.runFrom, Turing.MultiTapeTM.runFrom_succ_eq_step, Turing.MultiTapeTM.runFrom_succ_eq_step', Turing.MultiTapeTM.runFrom_zero, Turing.MultiTapeTM.step, Turing.moveInputPos_zero}
\difficulty{5}
Suppose the candidate word $s$ decomposes as
$\mathit{pre} \mathbin{+\!\!+} \mathit{rest}$, and let $\mathit{out}$ be a word and
$p$ an input-head position. After exactly $|\mathit{rest}| + 1$ steps from the
candidate-phase configuration of the assembly machine with input head at $p$,
candidate head at position $|\mathit{pre}|$, and output
$\mathit{out} \mathbin{+\!\!+} \mathit{pre}$, the machine is halted with the
candidate tape unchanged, its head at position $|s|$, the input head still at $p$,
and output $\mathit{out} \mathbin{+\!\!+} s$; an empty remainder takes the single
final step.
\end{lemma}

\begin{lemma}[Assembly emits the padded input]
\lean{Complexity.enumCont_concat_run}
\label{Complexity.enumCont_concat_run}
\leanok
\uses{Complexity.enumCont_concatCfg, Complexity.enumCont_concatTM, Complexity.enumCont_concat_candidate, Complexity.enumCont_concat_native, Turing.Cfg.ofWords, Turing.MultiTapeTM.runFrom, Turing.MultiTapeTM.runFrom_add}
\difficulty{3}
For all binary words $x$ and $s$, the assembly machine started in its input phase
on input $x$, with $s$ on its work tape and all heads at their starting positions,
is halted after exactly $|x| + |s| + 2$ steps with output $x \mathbin{+\!\!+} s$,
its input head at the right boundary, the candidate tape unchanged, and the
candidate head at position $|s|$.
\end{lemma}

\begin{lemma}[Buffered composition from a prepared first machine]
\lean{Complexity.enumCont_prepared_comp}
\label{Complexity.enumCont_prepared_comp}
\leanok
\uses{Complexity.enumCont_first_halt, Turing.Cfg, Turing.Cfg.init, Turing.FinTM, Turing.FinTM.ComputesInTime, Turing.FinTM.VirtualTag, Turing.FinTM.bufferedCompTM, Turing.FinTM.bufferedFirstCfg, Turing.FinTM.bufferedFirstCfg_rewind, Turing.FinTM.bufferedFirstCfg_run, Turing.FinTM.bufferedSecondCfg, Turing.FinTM.bufferedSecondCfg_run, Turing.FinTM.computesInTime_iff, Turing.MultiTapeTM.initCfg, Turing.MultiTapeTM.runFrom, Turing.MultiTapeTM.runFrom_add}
\difficulty{5}
Let $M_1, M_2$ be finite machines, and let $c_0$ be a configuration of $M_1$ on
input $x$ (not necessarily initial) from which $M_1$ halts within $T_1$ steps
with output $y$; suppose $M_2$ computes $z$ from input $y$ within $T_2$ steps.
Then the buffered composition of $M_1$ and $M_2$, started at its first-phase
configuration for $c_0$, halts within $T_1 + |y| + 2 + T_2$ steps with output $z$.
\end{lemma}

\begin{lemma}[Canonical candidate seam of a round]
\lean{Complexity.enumCont_round_seam}
\label{Complexity.enumCont_round_seam}
\leanok
\uses{Complexity.enumCont_concatTM, Turing.Cfg.ofWords, Turing.FinTM, Turing.FinTM.bufferedCompTM, Turing.FinTM.bufferedFirstCfg, Turing.FinTM.tapeBlocks, Turing.stateWord}
\difficulty{4}
For any verifier machine $MV$, words $x, s$, and assembly phase: the buffered
composition's first-phase configuration over the assembly machine prepared with
candidate $s$ equals the canonical configuration in which control sits in the
chosen assembly phase, tape zero carries $s$, and all other tapes (composition
buffer and verifier work) are blank.
\end{lemma}

\begin{lemma}[A prepared verifier call]
\lean{Complexity.enumCont_verifier_call}
\label{Complexity.enumCont_verifier_call}
\leanok
\uses{Complexity.enumCont_concatTM, Complexity.enumCont_concat_run, Complexity.enumCont_prepared_comp, Complexity.enumCont_round_seam, Turing.Cfg.ofWords, Turing.FinTM, Turing.FinTM.DecidesInTime, Turing.FinTM.bufferedCompTM, Turing.MultiTapeTM.indicator, Turing.MultiTapeTM.runFrom, Turing.stateWord}
\difficulty{3}
If a machine $MV$ decides a language $V$ within time $T_v$, then from the
canonical candidate seam with candidate $s$ on input $x$, the buffered
composition of the assembly machine with $MV$ halts within
$T_v(|x| + |s|) + 2(|x| + |s|) + 4$ steps, with single-bit output the indicator
of $x \mathbin{+\!\!+} s \in V$.
\end{lemma}

\begin{definition}[Halt-redirecting action wrapper]
\lean{Complexity.enumCont_returnAction}
\label{Complexity.enumCont_returnAction}
\leanok
\uses{Turing.Action}
Relabel the live states of an action along an embedding and redirect its halt to
a designated live return state, keeping all tape effects and the output
unchanged.
\end{definition}

\begin{definition}[Halt-redirecting configuration map]
\lean{Complexity.enumCont_returnCfg}
\label{Complexity.enumCont_returnCfg}
\leanok
\uses{Turing.Cfg}
The configuration correspondence matching the halt redirection: all fields are
preserved, and the control state is the embedded live state, or the designated
return state when the source has halted.
\end{definition}

\begin{lemma}[Simulation through the first halt via redirection]
\lean{Complexity.enumCont_return_run}
\label{Complexity.enumCont_return_run}
\leanok
\uses{Complexity.enumCont_returnAction, Complexity.enumCont_returnCfg, Turing.Cfg, Turing.Cfg.Halted, Turing.MultiTapeTM, Turing.MultiTapeTM.runFrom, Turing.MultiTapeTM.runFrom_succ_eq_step', Turing.MultiTapeTM.step}
\difficulty{4}
Suppose a host machine's transitions on embedded states are the halt-redirected
transitions of a source machine, and the source is live at every time before $t$
when run from $c_0$. Then the host's $t$-step run from the redirected
configuration of $c_0$ equals the redirected configuration of the source's
time-$t$ configuration; in particular at the first halt the host lands in the
return state with the source's final writes and output.
\end{lemma}

\begin{definition}[Padding an action to a larger tape bank]
\lean{Complexity.enumCont_padAction}
\label{Complexity.enumCont_padAction}
\leanok
\uses{Turing.Action}
Given a map $\iota$ from source control states to host control states, the padding
of an action $a$ of a $k$-tape machine to $K$ tapes is the action that performs
$a$'s input move, output, writes and moves on the first $k$ work tapes, leaves
every further tape unwritten and stationary, and replaces $a$'s next control state
by its image under $\iota$ (halting stays halting).
\end{definition}

\begin{definition}[Padding a configuration to a larger tape bank]
\lean{Complexity.enumCont_padCfg}
\label{Complexity.enumCont_padCfg}
\leanok
\uses{Turing.Cfg}
Pad a configuration of a $k$-tape machine to $K$ tapes by adding blank stationary
high tapes, embedding the control state along a given map.
\end{definition}

\begin{lemma}[Padding commutes with one action]
\lean{Complexity.enumCont_pad_apply}
\label{Complexity.enumCont_pad_apply}
\leanok
\uses{Complexity.enumCont_padAction, Complexity.enumCont_padCfg, Turing.Action, Turing.Action.apply, Turing.Cfg}
\difficulty{3}
Applying the padded version of an action to the padded version of a configuration
equals the padding of the result of applying the original action; the added high
tapes remain blank.
\end{lemma}

\begin{lemma}[Running a machine on an initial tape block of a host]
\lean{Complexity.enumCont_pad_run}
\label{Complexity.enumCont_pad_run}
\leanok
\uses{Complexity.enumCont_padAction, Complexity.enumCont_padCfg, Complexity.enumCont_pad_apply, Complexity.enumCont_sources, Turing.Cfg, Turing.Cfg.inputSymbol, Turing.Cfg.workTapeSymbols, Turing.MultiTapeTM, Turing.MultiTapeTM.runFrom, Turing.MultiTapeTM.runFrom_comm_of_step, Turing.MultiTapeTM.step}
\difficulty{4}
Let $k \le K$ and suppose a host machine's transitions on embedded states are the
padded transitions of a $k$-tape source machine, reading only the first $k$
tapes. Then for every configuration $c$ of the source and every time $t$, the
host's run from the padded configuration of $c$ equals the padding of the
source's run.
\end{lemma}

\begin{definition}[Three source routines on a shared tape bank]
\lean{Complexity.enumCont_sources}
\label{Complexity.enumCont_sources}
\leanok
\uses{Complexity.enumCont_padAction, Turing.FinTM}
The disjoint union of three finite machines on a common bank whose size is the
sum of their tape counts: each branch runs its machine, padded, on an initial
block of tapes. The machine starts in the third routine, and control can enter
either other routine without disturbing tape zero.
\end{definition}

\begin{definition}[Two-ended candidate/capture action]
\lean{Complexity.enumCont_endsAction}
\label{Complexity.enumCont_endsAction}
\leanok
\uses{Turing.Action}
Given optional writes $a$ and $b$, a displacement $d$, and a next control state,
the two-ended action is the action on $L + 1$ work tapes that performs write $a$ on
tape zero and write $b$ on the last tape, moves both of these heads by $d$, leaves
every intervening tape unwritten and stationary, keeps the input head fixed, and
emits nothing.
\end{definition}

\begin{definition}[The enumerator body]
\lean{Complexity.enumCont_bodyTM}
\label{Complexity.enumCont_bodyTM}
\leanok
\uses{Complexity.enumCont_cleanTM, Complexity.enumCont_endsAction, Complexity.enumCont_logTM, Complexity.enumCont_returnAction, Turing.FinTM, Turing.FinTM.controlAction}
Let $B$ be a finite source machine with designated verification and increment
entry states. The enumerator body built from $B$ is the machine that issues clean
calls of $B$ (with halts redirected to live administrative states) for
initialization, verification, and increment, coordinated by seven administrative
states: an anchor, a seed-copy phase, an unused reserve, a verdict test, an
increment check, a result copy, and a synchronized rewind. Its stopped variant
agrees with the active body except that the anchor is a silent self-loop.
\end{definition}

\begin{lemma}[A prepared increment call]
\lean{Complexity.enumCont_increment_call}
\label{Complexity.enumCont_increment_call}
\leanok
\uses{Complexity.enumCont_concatCfg, Complexity.enumCont_concatTM, Complexity.enumCont_concat_candidate, Complexity.enumCont_prepared_comp, Complexity.enumCont_round_seam, Turing.Cfg.ofWords, Turing.FinTM, Turing.FinTM.ComputesFunInTime, Turing.FinTM.bufferedCompTM, Turing.MultiTapeTM.runFrom, Turing.incFixed, Turing.stateWord}
\difficulty{4}
Suppose a machine $I$ computes, within time $j(n+1)$ on inputs of length $n$, the
output word of the fixed-width binary increment (the successor word, or the empty
word on overflow). Then from the candidate seam in the candidate phase with
candidate $s$ on input $x$, the buffered composition of the assembly machine
with $I$ halts within $(j+3)(|s|+1)$ steps and outputs the increment result of
$s$, which is empty exactly on overflow.
\end{lemma}

\begin{lemma}[Padding preserves the candidate seam]
\lean{Complexity.enumCont_pad_words}
\label{Complexity.enumCont_pad_words}
\leanok
\uses{Complexity.enumCont_padCfg, Turing.Cfg.ofWords, Turing.stateWord}
\difficulty{3}
For $0 < k \le K$: padding the canonical configuration whose only nonblank tape
is tape zero carrying a word $s$ from a $k$-tape bank to a $K$-tape bank yields
the canonical configuration with $s$ on tape zero of the larger bank, all added
tapes blank and parked at zero.
\end{lemma}

\begin{lemma}[Padding the initial configuration]
\lean{Complexity.enumCont_pad_init}
\label{Complexity.enumCont_pad_init}
\leanok
\uses{Complexity.enumCont_padCfg, Turing.Cfg.init, Turing.Cfg.ofWords, Turing.stateWord}
\difficulty{2}
Padding the initial (all-blank-work) configuration of a $k$-tape machine yields
the canonical configuration with an empty tape-zero word on the larger bank;
this holds even when the source machine has no work tapes.
\end{lemma}

\begin{lemma}[Overwriting the next candidate cell]
\lean{Complexity.enumCont_overwrite}
\label{Complexity.enumCont_overwrite}
\leanok
\uses{Turing.FinTM.bufferTape}
\difficulty{4}
Writing a bit at position $|\mathit{pre}|$ of the buffer tape holding
$\mathit{pre} \mathbin{+\!\!+} \mathit{old}$ yields the buffer tape holding
$\mathit{pre}$, then that bit, then the tail of $\mathit{old}$ with its first
cell dropped; this includes the case of an empty $\mathit{old}$.
\end{lemma}

\begin{lemma}[Clearing the next captured cell]
\lean{Complexity.enumCont_sparse_clear}
\label{Complexity.enumCont_sparse_clear}
\leanok
\uses{Complexity.enumCont_sparse}
\difficulty{4}
For every $p$, every bit $b$, and every boolean word $\mathit{rest}$, blanking
position $p$ of the sparse tape consisting of $p$ blank cells followed by the
symbols of the word $b$ then $\mathit{rest}$ yields the sparse tape of $p + 1$ blank
cells followed by the symbols of $\mathit{rest}$.
\end{lemma}

\begin{definition}[Two-ended administrative configurations]
\lean{Complexity.enumCont_pairCfg}
\label{Complexity.enumCont_pairCfg}
\leanok
\uses{Turing.Cfg}
A configuration on $L + 1$ work tapes for the administrative copy and rewind
scans: tape zero and the last tape carry given tape functions with a common head
position, every other tape is blank with its head at zero, the input head is at
cell one, and the output is empty.
\end{definition}

\begin{lemma}[Effect of a two-ended action]
\lean{Complexity.enumCont_ends_apply}
\label{Complexity.enumCont_ends_apply}
\leanok
\uses{Complexity.enumCont_endsAction, Complexity.enumCont_pairCfg, Turing.Action.apply, Turing.moveInputPos_zero}
\difficulty{4}
Let $L > 0$, and consider a two-ended configuration on $L + 1$ work tapes whose
tape zero holds $u$ and whose last tape holds $v$, with common head position $h$.
Applying the two-ended action with optional writes $a$, $b$, displacement $d$ and
next state $q'$ yields the two-ended configuration in state $q'$ whose tape zero is
$u$ with cell $h$ overwritten by $a$ (if $a$ writes), whose last tape is $v$ with
cell $h$ overwritten by $b$ (if $b$ writes), and whose common head position is
$h + d$; the input head and the empty output are unchanged.
\end{lemma}

\begin{lemma}[A list of blanks induces the blank tape]
\lean{Complexity.enumCont_sparse_blanks}
\label{Complexity.enumCont_sparse_blanks}
\leanok
\uses{Complexity.enumCont_sparse}
\difficulty{2}
The sparse tape induced by a list of $n$ blank entries is blank at every
position.
\end{lemma}

\begin{lemma}[Lifted words induce buffer tapes]
\lean{Complexity.enumCont_sparse_some}
\label{Complexity.enumCont_sparse_some}
\leanok
\uses{Complexity.enumCont_sparse, Turing.FinTM.bufferTape}
\difficulty{2}
The sparse tape induced by lifting every symbol of a boolean word into a
nonblank cell equals the standard buffer tape of that word.
\end{lemma}

\begin{lemma}[The copy scan replaces the captured suffix]
\lean{Complexity.enumCont_copy_scan}
\label{Complexity.enumCont_copy_scan}
\leanok
\uses{Complexity.enumCont_endsAction, Complexity.enumCont_ends_apply, Complexity.enumCont_overwrite, Complexity.enumCont_pairCfg, Complexity.enumCont_sparse, Complexity.enumCont_sparse_blanks, Complexity.enumCont_sparse_clear, Turing.Action.apply, Turing.Cfg.workTapeSymbols, Turing.FinTM.bufferTape, Turing.MultiTapeTM, Turing.MultiTapeTM.runFrom, Turing.MultiTapeTM.runFrom_succ_eq_step, Turing.MultiTapeTM.runFrom_succ_eq_step', Turing.MultiTapeTM.runFrom_zero, Turing.MultiTapeTM.step}
\difficulty{5}
Consider a machine with at least two work tapes whose copy state, at a two-ended
configuration, overwrites the tape-zero cell with $f$ applied to the captured
symbol, blanks the captured cell, and advances both heads, exiting left into a
rewind state at the captured right blank. Starting with tape zero holding
$\mathit{pre} \mathbin{+\!\!+} \mathit{old}$, the capture tape holding
$\mathit{rest}$ beyond $|\mathit{pre}|$ already-erased cells, and both heads at
$|\mathit{pre}|$: after $|\mathit{rest}| + 1$ steps the machine is in the rewind
state, tape zero holds $\mathit{pre}$, then $f$ mapped over $\mathit{rest}$, then
the remainder of $\mathit{old}$ beyond position $|\mathit{rest}|$, the capture
tape is blank, and both heads stand at $|\mathit{pre}| + |\mathit{rest}| - 1$.
\end{lemma}

\begin{lemma}[Synchronized rewind of the endpoint heads]
\lean{Complexity.enumCont_pair_rewind}
\label{Complexity.enumCont_pair_rewind}
\leanok
\uses{Complexity.enumCont_endsAction, Complexity.enumCont_ends_apply, Complexity.enumCont_pairCfg, Turing.Action.apply, Turing.Cfg.workTapeSymbols, Turing.FinTM.bufferTape, Turing.MultiTapeTM, Turing.MultiTapeTM.runFrom, Turing.MultiTapeTM.runFrom_succ_eq_step, Turing.MultiTapeTM.runFrom_succ_eq_step', Turing.MultiTapeTM.runFrom_zero, Turing.MultiTapeTM.step}
\difficulty{5}
Consider a machine whose rewind state moves both endpoint heads left while the
tape-zero cell is nonblank and exits right into an arrival state at the left
blank. For every word $w$ on tape zero and every $j \le |w|$: from common head
position $j - 1$, exactly $j + 1$ steps reach the arrival state with both heads
at zero, preserving $w$ and the blank capture tape.
\end{lemma}

\begin{lemma}[Complete copy and rewind]
\lean{Complexity.enumCont_copy_complete}
\label{Complexity.enumCont_copy_complete}
\leanok
\uses{Complexity.enumCont_copy_scan, Complexity.enumCont_endsAction, Complexity.enumCont_pairCfg, Complexity.enumCont_pair_rewind, Complexity.enumCont_sparse_some, Turing.FinTM.bufferTape, Turing.MultiTapeTM, Turing.MultiTapeTM.runFrom, Turing.MultiTapeTM.runFrom_add}
\difficulty{3}
For a machine with copy and rewind states as in the two preceding scans: if the
captured word $v$ is at least as long as the old tape-zero word, then from a
two-ended configuration with both heads at zero, exactly $2|v| + 2$ steps replace
tape zero entirely by $f$ mapped over $v$, blank the capture tape, and return
both heads to zero in the arrival state.
\end{lemma}

\begin{lemma}[Logged seam with empty history]
\lean{Complexity.enumCont_log_words}
\label{Complexity.enumCont_log_words}
\leanok
\uses{Complexity.enumCont_logCfg, Complexity.enumCont_logTM, Complexity.enumCont_sparse_blanks, Turing.Cfg.ofWords, Turing.FinTM, Turing.FinTM.tapeBlocks, Turing.stateWord}
\difficulty{4}
For a machine $B$ with at least one work tape: the logged configuration, with
empty history, of the canonical seam carrying a word $s$ on tape zero equals the
canonical seam of the logging machine's larger tape bank, the added tapes being
blank and stationary.
\end{lemma}

\begin{lemma}[Appending an output tape to a seam]
\lean{Complexity.enumCont_extend_words}
\label{Complexity.enumCont_extend_words}
\leanok
\uses{Turing.FinTM.bufferTape, Turing.stateWord}
\difficulty{3}
Appending one tape holding a word $y$ to a bank of $L$ tapes whose only nonblank
tape is tape zero holding $s$ yields exactly the two-ended tape layout: tape
zero $s$, last tape $y$, all intervening tapes blank.
\end{lemma}

\begin{lemma}[Clean entry at a prepared seam]
\lean{Complexity.enumCont_clean_entry_words}
\label{Complexity.enumCont_clean_entry_words}
\leanok
\uses{Complexity.enumCont_cleanTM, Complexity.enumCont_extend_words, Complexity.enumCont_logCfg, Complexity.enumCont_log_words, Complexity.enumCont_pairCfg, Turing.Cfg.ofWords, Turing.FinTM, Turing.FinTM.bufferTape, Turing.FinTM.bufferTape_nil, Turing.captureCfg, Turing.stateWord}
\difficulty{3}
For a machine $B$ with at least one work tape: the clean wrapper's captured entry
configuration at the prepared seam carrying $s$ on tape zero equals the
two-ended administrative configuration with tape zero $s$, blank capture, and
both endpoint heads at zero.
\end{lemma}

\begin{lemma}[Clean exit at a prepared seam]
\lean{Complexity.enumCont_clean_exit_words}
\label{Complexity.enumCont_clean_exit_words}
\leanok
\uses{Complexity.enumCont_cleanCfg, Complexity.enumCont_cleanTM, Complexity.enumCont_extend_words, Complexity.enumCont_logTM, Complexity.enumCont_log_words, Complexity.enumCont_pairCfg, Turing.Cfg.ofWords, Turing.FinTM, Turing.FinTM.bufferTape, Turing.stateWord}
\difficulty{3}
For a machine $B$ with at least one work tape: the clean wrapper's completed
administrative configuration over the prepared seam carrying $s$, with retained
output $y$, equals the two-ended configuration with tape zero $s$, capture tape
$y$, and both endpoint heads at zero.
\end{lemma}

\begin{lemma}[A clean call inside the body]
\lean{Complexity.enumCont_body_call}
\label{Complexity.enumCont_body_call}
\leanok
\uses{Complexity.enumCont_bodyTM, Complexity.enumCont_cleanTM, Complexity.enumCont_clean_complete, Complexity.enumCont_clean_entry_words, Complexity.enumCont_clean_exit_words, Complexity.enumCont_first_halt, Complexity.enumCont_pairCfg, Complexity.enumCont_return_run, Turing.Cfg.ofWords, Turing.FinTM, Turing.FinTM.bufferTape, Turing.MultiTapeTM.runFrom, Turing.stateWord}
\difficulty{4}
Let $B$ be a source machine with at least one work tape that, started in state
$q$ at the seam carrying candidate $s$ on input $x$, halts within $T$ steps with
output $y$. Then the enumerator body, from the two-ended configuration embedding
the clean call of $q$ in a given mode, reaches within $3T + |x| + |y| + 8$ steps
the two-ended configuration in that mode's administrative continuation state,
with candidate $s$ preserved, $y$ on the capture tape, all scratch blank, both
endpoint heads at zero, and empty physical output.
\end{lemma}

\begin{lemma}[Two-ended configurations with blank capture are seams]
\lean{Complexity.enumCont_pair_words}
\label{Complexity.enumCont_pair_words}
\leanok
\uses{Complexity.enumCont_pairCfg, Turing.Cfg.ofWords, Turing.FinTM.bufferTape, Turing.stateWord}
\difficulty{3}
The two-ended administrative configuration with tape zero holding $s$, blank
capture tape, and both endpoint heads at zero equals the canonical seam
configuration carrying $s$ on tape zero.
\end{lemma}

\begin{lemma}[Initialization of the enumerator body]
\lean{Complexity.enumCont_body_start}
\label{Complexity.enumCont_body_start}
\leanok
\uses{Complexity.enumCont_bodyTM, Complexity.enumCont_body_call, Complexity.enumCont_cleanTM, Complexity.enumCont_copy_complete, Complexity.enumCont_logTM, Complexity.enumCont_pairCfg, Complexity.enumCont_pair_words, Turing.Cfg.init, Turing.Cfg.ofWords, Turing.FinTM, Turing.FinTM.bufferTape, Turing.MultiTapeTM.initCfg, Turing.MultiTapeTM.runFrom, Turing.MultiTapeTM.runFrom_add, Turing.stateWord}
\difficulty{4}
Let $B$ be a source machine with at least one work tape that, started at the
empty seam on input $x$, halts within $T$ steps emitting $w$ consecutive true
marks. Then within $3T + |x| + 3w + 10$ steps from its initial configuration on
$x$, the enumerator body reaches the anchor state over the canonical seam
carrying the all-false candidate of width $w$.
\end{lemma}

\begin{lemma}[Departure from the anchor]
\lean{Complexity.enumCont_body_depart}
\label{Complexity.enumCont_body_depart}
\leanok
\uses{Complexity.enumCont_bodyTM, Complexity.enumCont_pairCfg, Complexity.enumCont_pair_words, Turing.Action.apply, Turing.Cfg.ofWords, Turing.FinTM, Turing.FinTM.bufferTape, Turing.FinTM.controlAction, Turing.FinTM.controlAction_apply, Turing.MultiTapeTM.step, Turing.moveInputPos_zero, Turing.stateWord}
\difficulty{3}
For every input $x$ and candidate word $s$, one step of the active enumerator body
from the anchor state over the canonical seam carrying $s$ on tape zero yields the
two-ended configuration that enters the verification call (the source's verifier
entry state in verification mode), with tape zero $s$, a blank capture tape, and
both endpoint heads at zero; the step writes nothing, moves no head, and emits
nothing.
\end{lemma}

\begin{lemma}[Acceptance emits the single true bit]
\lean{Complexity.enumCont_body_accept}
\label{Complexity.enumCont_body_accept}
\leanok
\uses{Complexity.enumCont_bodyTM, Complexity.enumCont_logTM, Complexity.enumCont_pairCfg, Turing.Action.apply, Turing.Cfg.workTapeSymbols, Turing.FinTM, Turing.FinTM.bufferTape, Turing.MultiTapeTM.step}
\difficulty{2}
For every candidate word $s$, after one step of the enumerator body (active or
stopped) from the two-ended configuration in the verdict state with tape zero $s$
and the captured word consisting of a single true bit, the machine is halted and
its output is exactly the single bit true.
\end{lemma}

\begin{lemma}[Rejection erases the verdict]
\lean{Complexity.enumCont_body_reject}
\label{Complexity.enumCont_body_reject}
\leanok
\uses{Complexity.enumCont_bodyTM, Complexity.enumCont_logTM, Complexity.enumCont_pairCfg, Turing.Action.apply, Turing.Cfg.workTapeSymbols, Turing.FinTM, Turing.FinTM.bufferTape, Turing.FinTM.bufferTape_append, Turing.MultiTapeTM.step, Turing.moveInputPos_zero}
\difficulty{4}
For every candidate word $s$, one step of the enumerator body (active or stopped)
from the two-ended configuration in the verdict state with tape zero $s$ and a
captured single false bit yields the two-ended configuration that enters the
increment call (the source's increment entry state in increment mode), with tape
zero $s$, the captured verdict erased so that the capture tape is blank, and both
endpoint heads at zero.
\end{lemma}

\begin{lemma}[Case analysis of the increment result]
\lean{Complexity.enumCont_increment_cases}
\label{Complexity.enumCont_increment_cases}
\leanok
\uses{Complexity.enumCont_step_length, Turing.incFixed}
\difficulty{3}
For every boolean word $s$, exactly one of the following holds for the
fixed-width increment's output word: it is empty (overflow) and the stalled step
leaves $s$ unchanged; or it is nonempty, has length exactly $|s|$, and the
stalled step equals it.
\end{lemma}

\begin{lemma}[The increment-result check]
\lean{Complexity.enumCont_body_check}
\label{Complexity.enumCont_body_check}
\leanok
\uses{Complexity.enumCont_bodyTM, Complexity.enumCont_logTM, Complexity.enumCont_pairCfg, Turing.Action.apply, Turing.Cfg.workTapeSymbols, Turing.FinTM, Turing.FinTM.bufferTape, Turing.FinTM.controlAction_apply, Turing.MultiTapeTM.step, Turing.moveInputPos_zero}
\difficulty{4}
One step of the enumerator body from the increment-check state, at the two-ended
configuration with candidate $s$ and captured increment result $v$, preserves
all tapes and transfers control to the anchor when $v$ is empty (overflow) and
to the result-copy state otherwise.
\end{lemma}

\begin{lemma}[One complete candidate update]
\lean{Complexity.enumCont_body_increment}
\label{Complexity.enumCont_body_increment}
\leanok
\uses{Complexity.enumCont_bodyTM, Complexity.enumCont_body_agree, Complexity.enumCont_body_call, Complexity.enumCont_body_check, Complexity.enumCont_copy_complete, Complexity.enumCont_increment_cases, Complexity.enumCont_logTM, Complexity.enumCont_pairCfg, Complexity.enumCont_pair_words, Turing.Cfg.ofWords, Turing.FinTM, Turing.FinTM.bufferTape, Turing.FinTM.bufferTape_nil, Turing.MultiTapeTM.runFrom, Turing.MultiTapeTM.runFrom_add, Turing.MultiTapeTM.runFrom_succ_eq_step', Turing.incFixed, Turing.stateWord}
\difficulty{5}
Let $B$ be a source machine with at least one work tape that, started in the
increment entry state at the seam with candidate $s$, halts within $T$ steps
emitting the fixed-width increment's output word. Then the enumerator body, from
the two-ended configuration embedding the increment call, reaches within
$3T + |x| + 3|s| + 11$ steps the anchor over the seam carrying the stalled
increment of $s$: the successor word of the same width, or $s$ itself on
overflow.
\end{lemma}

\begin{lemma}[One raw round after departure]
\lean{Complexity.enumCont_body_round_raw}
\label{Complexity.enumCont_body_round_raw}
\leanok
\uses{Complexity.enumCont_bodyTM, Complexity.enumCont_body_accept, Complexity.enumCont_body_call, Complexity.enumCont_body_increment, Complexity.enumCont_body_reject, Complexity.enumCont_pairCfg, Turing.Cfg.ofWords, Turing.FinTM, Turing.FinTM.bufferTape, Turing.MultiTapeTM.runFrom, Turing.MultiTapeTM.runFrom_add, Turing.MultiTapeTM.runFrom_succ_eq_step', Turing.incFixed, Turing.stateWord}
\difficulty{4}
Let $B$ be a source machine with at least one work tape whose verification call
on candidate $s$ (from the verifier entry state at the seam) halts within $T_v$
steps with single-bit output $b$, and whose increment call halts within $T_i$
steps emitting the fixed-width increment output. Then from the configuration
just after departure, within $3T_v + 3T_i + 2|x| + 3|s| + 21$ steps the
enumerator body either halts with output true (when $b$ is true) or reaches the
anchor over the seam carrying the stalled increment of $s$ (when $b$ is false).
\end{lemma}

\begin{lemma}[A silent self-loop absorbs]
\lean{Complexity.enumCont_absorb}
\label{Complexity.enumCont_absorb}
\leanok
\uses{Turing.Cfg, Turing.FinTM.controlAction, Turing.FinTM.controlAction_apply, Turing.MultiTapeTM, Turing.MultiTapeTM.runFrom, Turing.MultiTapeTM.runFrom_succ_eq_step', Turing.MultiTapeTM.step, Turing.moveInputPos_zero}
\difficulty{3}
If a control state's transition is the silent stationary self-loop, then from
any configuration in that state, a run of any length leaves the configuration
unchanged.
\end{lemma}

\begin{lemma}[Runs agree before the first anchor visit]
\lean{Complexity.enumCont_agree_run}
\label{Complexity.enumCont_agree_run}
\leanok
\uses{Turing.Cfg, Turing.MultiTapeTM, Turing.MultiTapeTM.runFrom, Turing.MultiTapeTM.runFrom_succ_eq_step', Turing.MultiTapeTM.step}
\difficulty{4}
If two transition tables agree at every control state other than a designated
anchor state, then their runs from a common configuration coincide for every
duration during which the second machine's run has not yet visited the anchor.
\end{lemma}

\begin{lemma}[First visit to a stopped anchor]
\lean{Complexity.enumCont_first_anchor}
\label{Complexity.enumCont_first_anchor}
\leanok
\uses{Complexity.enumCont_absorb, Complexity.enumCont_agree_run, Turing.Cfg, Turing.FinTM.controlAction, Turing.MultiTapeTM, Turing.MultiTapeTM.runFrom, Turing.MultiTapeTM.runFrom_add}
\difficulty{4}
Let two transition tables agree away from an anchor state, which the second
(stopped) table turns into a silent self-loop. If the stopped run from $c$
stands at the anchor at time $T$, then there is a time $t \le T$ such that the
active run never visits the anchor before $t$ and its time-$t$ configuration
equals the stopped run's time-$T$ configuration.
\end{lemma}

\begin{lemma}[Halted stopped runs transfer unchanged]
\lean{Complexity.enumCont_halt_transfer}
\label{Complexity.enumCont_halt_transfer}
\leanok
\uses{Complexity.enumCont_absorb, Complexity.enumCont_agree_run, Turing.Cfg, Turing.FinTM.controlAction, Turing.MultiTapeTM, Turing.MultiTapeTM.runFrom, Turing.MultiTapeTM.runFrom_add}
\difficulty{4}
Let two transition tables agree away from an anchor state, which the second
(stopped) table turns into a silent self-loop. If the stopped run from $c$ is
halted at time $T$, then the active run from $c$ equals the stopped run at time
$T$, and the active run never visits the anchor at any time up to $T$.
\end{lemma}

\begin{lemma}[Active and stopped bodies agree off the anchor]
\lean{Complexity.enumCont_body_agree}
\label{Complexity.enumCont_body_agree}
\leanok
\uses{Complexity.enumCont_bodyTM, Turing.FinTM}
\difficulty{3}
The active and stopped variants of the enumerator body have identical transitions
at every control state other than the anchor; their tape counts and finite state
sets coincide.
\end{lemma}

\begin{lemma}[Guarded initialization of the body]
\lean{Complexity.enumCont_body_start_guarded}
\label{Complexity.enumCont_body_start_guarded}
\leanok
\uses{Complexity.enumCont_bodyTM, Complexity.enumCont_body_agree, Complexity.enumCont_body_start, Complexity.enumCont_first_anchor, Turing.Cfg.ofWords, Turing.FinTM, Turing.MultiTapeTM.initCfg, Turing.MultiTapeTM.runFrom, Turing.stateWord}
\difficulty{3}
Let $B$ be a source machine with at least one work tape that, from the empty
seam on input $x$, halts within $T$ steps emitting $w$ true marks. Then within
$3T + |x| + 3w + 10$ steps the active enumerator body reaches the anchor over
the all-false width-$w$ candidate seam, without visiting the anchor at any
earlier time.
\end{lemma}

\begin{lemma}[Guarded round of the body]
\lean{Complexity.enumCont_body_round_guarded}
\label{Complexity.enumCont_body_round_guarded}
\leanok
\uses{Complexity.enumCont_bodyTM, Complexity.enumCont_body_agree, Complexity.enumCont_body_depart, Complexity.enumCont_body_round_raw, Complexity.enumCont_first_anchor, Complexity.enumCont_halt_transfer, Complexity.enumCont_pairCfg, Turing.Cfg.ofWords, Turing.FinTM, Turing.FinTM.bufferTape, Turing.MultiTapeTM.runFrom, Turing.MultiTapeTM.runFrom_succ_eq_step, Turing.incFixed, Turing.stateWord}
\difficulty{5}
Let $B$ be a source machine with at least one work tape whose verification call
on candidate $s$ halts within $T_v$ steps with single-bit verdict $b$, and whose
increment call halts within $T_i$ steps emitting the fixed-width increment
output. Then there is a time $t$ with
$0 < t \le 3T_v + 3T_i + 2|x| + 3|s| + 22$ such that the active enumerator body,
started at the anchor over the seam with candidate $s$, does not visit the
anchor at any strictly intermediate positive time, and at time $t$ either halts
with output true (when $b$ is true) or stands at the anchor over the seam
carrying the stalled increment of $s$ (when $b$ is false).
\end{lemma}

\begin{lemma}[Lifting a prepared call into a shared bank]
\lean{Complexity.enumCont_lift_call}
\label{Complexity.enumCont_lift_call}
\leanok
\uses{Complexity.enumCont_padAction, Complexity.enumCont_padCfg, Complexity.enumCont_pad_run, Complexity.enumCont_pad_words, Turing.Cfg.ofWords, Turing.FinTM, Turing.MultiTapeTM.runFrom, Turing.stateWord}
\difficulty{2}
Suppose a host machine simulates a source machine with at least one work tape by
padded transitions on an initial tape block. If the source, started in state $q$
at the seam carrying $s$, is halted at time $T$ with output $y$, then the host,
started in the embedded state at its own seam carrying $s$, is halted at time
$T$ with the same output $y$.
\end{lemma}

\begin{lemma}[Lifting an initial call into a shared bank]
\lean{Complexity.enumCont_lift_init}
\label{Complexity.enumCont_lift_init}
\leanok
\uses{Complexity.enumCont_padAction, Complexity.enumCont_padCfg, Complexity.enumCont_pad_init, Complexity.enumCont_pad_run, Turing.Cfg.ofWords, Turing.FinTM, Turing.FinTM.ComputesInTime, Turing.FinTM.computesInTime_iff, Turing.MultiTapeTM.initCfg, Turing.MultiTapeTM.runFrom, Turing.stateWord}
\difficulty{3}
Suppose a host machine simulates a source machine by padded transitions on an
initial tape block. If the source computes $y$ from input $x$ within $T$ steps,
then the host, started in the embedded initial state at the empty seam on $x$,
is halted at time $T$ with output $y$; this holds even for a source with no work
tapes.
\end{lemma}

\begin{lemma}[From body contracts to the enumerator contract]
\lean{Complexity.enumCont_from_body}
\label{Complexity.enumCont_from_body}
\leanok
\uses{Complexity.enumCont_fuel_bits, Complexity.enumCont_orbit, Complexity.enumCont_step_length, Complexity.enumWord, Turing.Cfg, Turing.Cfg.ofWords, Turing.FinTM, Turing.FinTM.ComputesFunInTime, Turing.FinTM.ComputesInTime.mono, Turing.FinTM.computesFunInTime_polyUnary, Turing.FinTM.exists_loopCfgTM, Turing.MultiTapeTM.indicator, Turing.MultiTapeTM.initCfg, Turing.MultiTapeTM.runFrom, Turing.incFixed, Turing.stateWord}
\difficulty{6}
Fix constants $C, c$, a budget function $G$, a language $V$ over the binary
alphabet, and a finite machine (the body) with a designated anchor state
satisfying: (i) from its initial configuration on any input $x$ it reaches,
within $G(|x|)$ steps and with no earlier anchor visit, the anchor over the seam
carrying the all-false candidate of width $w = C(|x|+1)^c$; and (ii) from the
anchor over any candidate $s$ of width $w$, within some positive time at most
$G(|x|)$ and with no interior anchor visit, it halts with output true when
$x \mathbin{+\!\!+} s \in V$ and otherwise returns to the anchor over the stalled
fixed-width increment of $s$. Then there exist a constant $b$ and a machine $E$
such that for every input $x$ there is a rank-indexed family of configurations
with: $E$ reaches the rank-$0$ configuration from its initial configuration
within $b(G(|x|) + (|x|+1)^{c+1} + 1)$ steps; the rank-$2^w$ configuration is
halted with single-bit output false; and from each rank $i < 2^w$, within the
same budget, $E$ either halts with output true (when $x$ concatenated with the
width-$w$ binary word of $i$ lies in $V$) or reaches rank $i + 1$.
\end{lemma}

\begin{theorem}[The enumerator's configuration contract]
\lean{Complexity.enumMachine_contracts}
\label{Complexity.enumMachine_contracts}
\leanok
\uses{Complexity.enumCont_bodyTM, Complexity.enumCont_body_round_guarded, Complexity.enumCont_body_start_guarded, Complexity.enumCont_concatTM, Complexity.enumCont_from_body, Complexity.enumCont_increment_call, Complexity.enumCont_lift_call, Complexity.enumCont_lift_init, Complexity.enumCont_sources, Complexity.enumCont_verifier_call, Complexity.enumWord, Turing.Cfg, Turing.FinTM, Turing.FinTM.DecidesInTime, Turing.FinTM.bufferedCompTM, Turing.FinTM.computesFunInTime_incFixed, Turing.FinTM.computesFunInTime_polyUnary, Turing.MultiTapeTM.indicator, Turing.MultiTapeTM.initCfg, Turing.MultiTapeTM.runFrom, Turing.incFixed}
\difficulty{6}
Fix constants $C, c$, a language $V$ over the binary alphabet, and a machine $MV$
deciding $V$ within time $T_v$. Then there exist a constant $b$ and one finite
machine $E$ such that for every input $x$ --- writing $n = |x|$ and
$w = C(n+1)^c$ --- there is a rank-indexed family of configurations of $E$ on
$x$ with: $E$ reaches the rank-$0$ configuration from its initial configuration
within $b\,(T_v(n+w) + (n+w+1)^{c+1})$ steps; the rank-$2^w$ configuration is
halted with single-bit output false; and from each rank $i < 2^w$, within the
same budget, $E$ either halts with output true when $x$ concatenated with the
width-$w$ binary word of $i$ lies in $V$, or reaches the rank-$(i+1)$
configuration otherwise.
\end{theorem}

\begin{theorem}[Brute-force enumeration decides the existential projection]
\lean{Complexity.exists_proj_decider}
\label{Complexity.exists_proj_decider}
\leanok
\uses{Complexity.enumAny, Complexity.enumAny_certificates, Complexity.enumLoop_run, Complexity.enumMachine_contracts, Complexity.enumWord, Turing.FinTM, Turing.FinTM.ComputesInTime, Turing.FinTM.ComputesInTime.mono, Turing.FinTM.DecidesInTime, Turing.FinTM.computesInTime_iff, Turing.MultiTapeTM.indicator, Turing.MultiTapeTM.runFrom_add}
\difficulty{5}
Let a machine $MV$ decide a language $V$ over the binary alphabet within time
$T_v$, and fix constants $C, c$. Then there exist a constant $b$ and a finite
machine deciding the existential projection
$\{\,x \mid \exists u,\ |u| = C(|x|+1)^c \wedge x \mathbin{+\!\!+} u \in V\,\}$
within time
\[
b \cdot 2^{C(n+1)^c} \cdot
  \Bigl(T_v\bigl(n + C(n+1)^c\bigr) + \bigl(n + C(n+1)^c + 1\bigr)^{c+1}\Bigr).
\]
\end{theorem}

\begin{theorem}[$\mathrm{NP} \subseteq \mathrm{EXP}$]
\lean{Complexity.NP_subset_EXP}
\label{Complexity.NP_subset_EXP}
\leanok
\uses{Complexity.EXP, Complexity.NP, Complexity.enumBudget_bound, Complexity.exists_proj_decider, Complexity.mem_P_iff, Turing.FinTM.ComputesInTime.mono}
\difficulty{4}
Every language in $\mathrm{NP}$ lies in $\mathrm{EXP}$: a language over the
binary alphabet with polynomial-length certificates and a polynomial-time
verifier is decidable in time $2^{n^e}$ for some constant $e$, up to a constant
factor.
\end{theorem}

\begin{lemma}[Binary digits of a power-of-two multiple]
\lean{Complexity.a3_bits_shift}
\label{Complexity.a3_bits_shift}
\leanok
\difficulty{3}
For every $C \neq 0$ and every $p$, the binary-digit list of $C \cdot 2^p$
consists of $p$ false bits followed by the binary digits of $C$.
\end{lemma}

\begin{definition}[The binary shift machine]
\lean{Complexity.a3ShiftTM}
\label{Complexity.a3ShiftTM}
\leanok
\uses{Turing.FinTM, Turing.FinTM.controlAction, Turing.FinTM.emitAction}
For a fixed boolean word $w$, a machine with no work tapes that emits one false
bit for each native input symbol and then emits $w$, halting afterwards.
\end{definition}

\begin{lemma}[Scanner invariant of the shift machine]
\lean{Complexity.a3_shift_scan}
\label{Complexity.a3_shift_scan}
\leanok
\uses{Complexity.a3ShiftTM, Turing.Action.apply, Turing.Cfg.workTapeSymbols, Turing.MultiTapeTM.initCfg, Turing.MultiTapeTM.runFrom, Turing.MultiTapeTM.runFrom_succ_eq_step', Turing.MultiTapeTM.step, Turing.inputSymbolInner, Turing.moveInputPos_pos_of_ne_right}
\difficulty{5}
For every $t \le |x|$: after $t$ steps on input $x$, the shift machine is still
in its scanning state, its input head stands at position $t + 1$, and its output
so far is exactly $t$ false bits.
\end{lemma}

\begin{lemma}[The shift machine computes its padded word]
\lean{Complexity.a3_shift_computes}
\label{Complexity.a3_shift_computes}
\leanok
\uses{Complexity.a3ShiftTM, Complexity.a3_shift_scan, Turing.Action.apply, Turing.Cfg.inputSymbol, Turing.FinTM.ComputesInTime, Turing.FinTM.computesInTime_iff, Turing.FinTM.controlAction, Turing.FinTM.emit_halts, Turing.MultiTapeTM.initCfg, Turing.MultiTapeTM.runFrom, Turing.MultiTapeTM.runFrom_add, Turing.MultiTapeTM.runFrom_succ_eq_step', Turing.MultiTapeTM.step}
\difficulty{5}
For every suffix word $w$ and every input $x$, the shift machine for $w$ halts on
input $x$ within $|x| + |w| + 2$ steps, and its output is $|x|$ false bits followed
by $w$.
\end{lemma}

\begin{lemma}[Monotone linear budget for the shift machine]
\lean{Complexity.a3_shift_timed}
\label{Complexity.a3_shift_timed}
\leanok
\uses{Complexity.a3ShiftTM, Complexity.a3_shift_computes, Turing.FinTM.ComputesFunInTime, Turing.FinTM.ComputesInTime.mono}
\difficulty{2}
For every word $w$, the shift machine for $w$ computes the function mapping each
input $x$ to $0^{|x|} \mathbin{+\!\!+} w$ within time $(|w| + 2)(n + 1)$ on inputs
of length $n$; this time bound is monotone and linear in $n$.
\end{lemma}

\begin{lemma}[Exponential padding widths are computable in binary]
\lean{Complexity.a3_exp_bits_timed}
\label{Complexity.a3_exp_bits_timed}
\leanok
\uses{Complexity.a3_bits_shift, Complexity.a3_shift_timed, Turing.FinTM, Turing.FinTM.ComputesFunInTime, Turing.FinTM.ComputesInTime.mono, Turing.FinTM.computesFunInTime_comp, Turing.FinTM.computesFunInTime_const, Turing.FinTM.computesFunInTime_polyUnary}
\difficulty{5}
For all constants $C, c$ there exist a finite machine and a constant $A$ such
that the machine computes, from each input $x$, the binary-digit list of
$C \cdot 2^{(|x|+1)^c}$ within time $A (|x| + 1)^{c+1}$.
\end{lemma}

\begin{definition}[Decider wrapper of an installed call]
\lean{Complexity.a3RunTM}
\label{Complexity.a3RunTM}
\leanok
\uses{Turing.Action.mapState, Turing.FinTM}
For a machine $C$ with at least one work tape and designated entry and exit
states: the wrapper machine that copies the native input onto tape zero, rewinds
both heads, executes one entry action of $C$, follows $C$ until its first return
to the exit state, and then emits the single symbol found on tape zero as its
decision.
\end{definition}

\begin{definition}[Load-phase configurations of the decider wrapper]
\lean{Complexity.a3LoadCfg}
\label{Complexity.a3LoadCfg}
\leanok
\uses{Turing.Cfg, Turing.FinTM, Turing.FinTM.bufferTape}
A load-phase configuration of the decider wrapper on input $x$ is determined by a
phase index $q \in \{0, 1, 2\}$, an input-head position $p$, a word $u$, and a head
position $h$: its control state is phase $q$, its input head is at $p$, tape zero
holds $u$ with its head at $h$, every other work tape is blank with its head at
zero, and the output is empty.
\end{definition}

\begin{lemma}[One copy transition]
\lean{Complexity.a3_copy_step}
\label{Complexity.a3_copy_step}
\leanok
\uses{Complexity.a3LoadCfg, Complexity.a3RunTM, Turing.Action.apply, Turing.Cfg.inputSymbol, Turing.FinTM, Turing.FinTM.bufferTape_append, Turing.MultiTapeTM.step, Turing.inputSymbolInner, Turing.moveInputPos, Turing.moveInputPos_pos_of_ne_right}
\difficulty{4}
For every $i < |x|$, one step of the decider wrapper in its copy phase, with the
first $i$ input bits already on tape zero, appends the $(i+1)$-st input bit to
tape zero and advances both the input head and the tape-zero head by one.
\end{lemma}

\begin{lemma}[The copy phase from blank tapes]
\lean{Complexity.a3_copy_run}
\label{Complexity.a3_copy_run}
\leanok
\uses{Complexity.a3LoadCfg, Complexity.a3RunTM, Complexity.a3_copy_step, Turing.Cfg.init, Turing.FinTM, Turing.FinTM.bufferTape, Turing.MultiTapeTM.initCfg, Turing.MultiTapeTM.runFrom, Turing.MultiTapeTM.runFrom_succ_eq_step'}
\difficulty{3}
For every $i \le |x|$, running the decider wrapper for $i$ steps from its
initial configuration on $x$ reaches the copy-phase configuration with the first
$i$ input bits on tape zero and both active heads just past them.
\end{lemma}

\begin{lemma}[The synchronized load rewind]
\lean{Complexity.a3_load_rewind}
\label{Complexity.a3_load_rewind}
\leanok
\uses{Complexity.a3LoadCfg, Complexity.a3RunTM, Turing.Action.apply, Turing.Cfg.ofWords, Turing.Cfg.workTapeSymbols, Turing.FinTM, Turing.FinTM.bufferTape, Turing.FinTM.bufferTape_left, Turing.FinTM.bufferTape_nat, Turing.FinTM.moveInputPos_neg_val, Turing.MultiTapeTM.runFrom, Turing.MultiTapeTM.runFrom_succ_eq_step, Turing.MultiTapeTM.runFrom_zero, Turing.MultiTapeTM.step, Turing.moveInputPos, Turing.moveInputPos_pos_of_ne_right, Turing.stateWord}
\difficulty{5}
For every $j \le |x|$: from the rewind-phase configuration of the decider
wrapper with the full copy of $x$ on tape zero and both active heads just left
of cell $j$, exactly $j + 1$ steps reach the dispatch seam --- the dispatch
state over $x$ on tape zero with all heads at the origin. This includes the
empty input.
\end{lemma}

\begin{lemma}[Linear-time arrival at the call seam]
\lean{Complexity.a3_run_start}
\label{Complexity.a3_run_start}
\leanok
\uses{Complexity.a3LoadCfg, Complexity.a3RunTM, Complexity.a3_copy_run, Complexity.a3_load_rewind, Turing.Action.apply, Turing.Cfg.inputSymbol, Turing.Cfg.ofWords, Turing.FinTM, Turing.FinTM.moveInputPos_neg_val, Turing.MultiTapeTM.initCfg, Turing.MultiTapeTM.runFrom, Turing.MultiTapeTM.runFrom_add, Turing.MultiTapeTM.runFrom_succ_eq_step', Turing.MultiTapeTM.step, Turing.moveInputPos, Turing.stateWord}
\difficulty{4}
Let $C$ be a finite machine with at least one work tape and designated entry and
exit states, and let $x$ be an input. After exactly $2|x| + 2$ steps from its
initial configuration on $x$, the decider wrapper of $C$ is in the dispatch seam
configuration: the dispatch state, $x$ on tape zero, every other work tape blank,
and all heads at their starting positions.
\end{lemma}

\begin{lemma}[Following the call until its first return]
\lean{Complexity.a3_run_guarded}
\label{Complexity.a3_run_guarded}
\leanok
\uses{Complexity.a3RunTM, Turing.Cfg, Turing.Cfg.mapState, Turing.FinTM, Turing.MultiTapeTM.runFrom, Turing.MultiTapeTM.runFrom_succ_eq_step', Turing.MultiTapeTM.step}
\difficulty{4}
As long as the inner machine's run from a configuration has not reached the
designated exit state, the decider wrapper's run from the state-renamed copy of
that configuration coincides, under the renaming, with the inner machine's run.
\end{lemma}

\begin{lemma}[Dispatch and first positive return]
\lean{Complexity.a3_run_call}
\label{Complexity.a3_run_call}
\leanok
\uses{Complexity.a3RunTM, Complexity.a3_run_guarded, Turing.Cfg.mapState, Turing.Cfg.ofWords, Turing.FinTM, Turing.MultiTapeTM.runFrom, Turing.MultiTapeTM.runFrom_succ_eq_step, Turing.MultiTapeTM.step, Turing.stateWord}
\difficulty{5}
Suppose the inner machine, started in its entry state at the seam carrying $x$
on tape zero, first reaches the exit state at a positive time $t$ (no visit at
any earlier positive time), in the canonical configuration carrying a word $y$
on tape zero. Then the decider wrapper, from its dispatch seam, reaches the
renamed exit configuration in exactly $t$ steps; this covers the case where the
entry and exit states coincide.
\end{lemma}

\begin{lemma}[Extracting the installed singleton verdict]
\lean{Complexity.a3_run_singleton}
\label{Complexity.a3_run_singleton}
\leanok
\uses{Complexity.a3RunTM, Complexity.a3_run_call, Complexity.a3_run_start, Turing.Action.apply, Turing.Cfg.mapState, Turing.Cfg.ofWords, Turing.Cfg.workTapeSymbols, Turing.FinTM, Turing.FinTM.ComputesInTime, Turing.FinTM.bufferTape, Turing.FinTM.computesInTime_iff, Turing.MultiTapeTM.runFrom, Turing.MultiTapeTM.runFrom_add, Turing.MultiTapeTM.runFrom_succ_eq_step', Turing.MultiTapeTM.step, Turing.stateWord}
\difficulty{3}
Suppose the inner machine, started in its entry state at the seam carrying $x$,
first returns to the exit state at a positive time $t$ with the singleton word
$[b]$ on tape zero. Then the decider wrapper computes the output $[b]$ from
input $x$ within $2|x| + 2 + t + 1$ steps, with no earlier emission.
\end{lemma}

\begin{lemma}[A clean installed call yields an ordinary decider]
\lean{Complexity.a3_decider_clean}
\label{Complexity.a3_decider_clean}
\leanok
\uses{Complexity.a3RunTM, Complexity.a3_run_singleton, Turing.FinTM, Turing.FinTM.ComputesInTime.mono, Turing.FinTM.DecidesInTime, Turing.FinTM.exists_installCallTM, Turing.MultiTapeTM.indicator}
\difficulty{4}
If a finite machine $M$ decides a language $L$ over the binary alphabet within
time $T$, then there exist a finite machine $R$ and a constant $B$ such that $R$
decides $L$ within time $B\,(T(n) + n + 2)$.
\end{lemma}

\begin{lemma}[Timed buffered composition at the actual intermediate word]
\lean{Complexity.a3_comp_at}
\label{Complexity.a3_comp_at}
\leanok
\uses{Turing.Cfg.init, Turing.FinTM, Turing.FinTM.ComputesInTime, Turing.FinTM.ComputesInTime.mono, Turing.FinTM.VirtualTag, Turing.FinTM.bufferedCompTM, Turing.FinTM.bufferedComp_start, Turing.FinTM.bufferedSecondCfg, Turing.FinTM.bufferedSecondCfg_run, Turing.FinTM.computesInTime_iff, Turing.MultiTapeTM.initCfg, Turing.MultiTapeTM.runFrom_add}
\difficulty{3}
If a machine $F$ computes $y$ from $x$ within $s$ steps and a machine $G$
computes $z$ from $y$ within $t$ steps, then the buffered composition of $F$ and
$G$ computes $z$ from $x$ within $s + |y| + 2 + t$ steps.
\end{lemma}

\begin{definition}[Nonempty-input guard]
\lean{Complexity.a3NonemptyTM}
\label{Complexity.a3NonemptyTM}
\leanok
\uses{Turing.Action.mapState, Turing.FinTM, Turing.FinTM.controlAction}
For a finite machine $M$, the nonempty-input guard of $M$ is the machine with the
same work tapes that, on empty input, emits a single rejecting (false) bit and
halts, and, on nonempty input, takes one silent step and then runs $M$ from $M$'s
initial state.
\end{definition}

\begin{lemma}[The guard preserves the source run]
\lean{Complexity.a3_nonempty_run}
\label{Complexity.a3_nonempty_run}
\leanok
\uses{Complexity.a3NonemptyTM, Turing.Cfg, Turing.Cfg.mapState, Turing.FinTM, Turing.MultiTapeTM.runFrom, Turing.MultiTapeTM.runFrom_comm_of_step, Turing.MultiTapeTM.step}
\difficulty{3}
After the nonempty guard has been passed, the guarded machine's run from any
state-renamed configuration of $M$ coincides, under the renaming, with $M$'s own
run.
\end{lemma}

\begin{lemma}[The guard rejects the empty input]
\lean{Complexity.a3_nonempty_nil}
\label{Complexity.a3_nonempty_nil}
\leanok
\uses{Complexity.a3NonemptyTM, Turing.Action.apply, Turing.Cfg.init, Turing.Cfg.inputSymbol, Turing.FinTM, Turing.FinTM.ComputesInTime, Turing.FinTM.computesInTime_iff, Turing.MultiTapeTM.initCfg, Turing.MultiTapeTM.runFrom, Turing.MultiTapeTM.step}
\difficulty{2}
For every finite machine $M$, the nonempty-input guard of $M$, run on the empty
input, is halted after one step with output exactly the single bit false.
\end{lemma}

\begin{lemma}[The guard passes nonempty inputs through]
\lean{Complexity.a3_nonempty_computes}
\label{Complexity.a3_nonempty_computes}
\leanok
\uses{Complexity.a3NonemptyTM, Complexity.a3_nonempty_run, Turing.Action.apply, Turing.Cfg.init, Turing.Cfg.inputSymbol, Turing.Cfg.mapState, Turing.FinTM, Turing.FinTM.ComputesInTime, Turing.FinTM.computesInTime_iff, Turing.FinTM.controlAction, Turing.MultiTapeTM.initCfg, Turing.MultiTapeTM.runFrom_succ_eq_step, Turing.MultiTapeTM.step, Turing.moveInputPos_zero}
\difficulty{3}
If $x$ is nonempty and $M$ computes $y$ from $x$ within $t$ steps, then the
guarded machine computes $y$ from $x$ within $t + 1$ steps.
\end{lemma}

\begin{lemma}[Strict monotonicity of exponential padding]
\lean{Complexity.a3_split_strictMono}
\label{Complexity.a3_split_strictMono}
\leanok
\difficulty{2}
For all constants $C, c$, the map $n \mapsto n + C \cdot 2^{(n+1)^c}$ on the
natural numbers is strictly increasing; in particular the exponential split of
a padded length, when it exists, is unique. This includes coefficient zero and
degree zero.
\end{lemma}

\begin{definition}[The exponential-split search]
\lean{Complexity.a3Split}
\label{Complexity.a3Split}
\leanok
Given constants $C, c$ and a length $m$: the least $n \le m$ satisfying
$n + C \cdot 2^{(n+1)^c} = m$, if one exists, and an explicit failure value
otherwise.
\end{definition}

\begin{lemma}[A successful split certifies its equation]
\lean{Complexity.a3_split_spec}
\label{Complexity.a3_split_spec}
\leanok
\uses{Complexity.a3Split}
\difficulty{2}
If the exponential-split search on $(C, c, m)$ returns $n$, then $n \le m$ and
$n + C \cdot 2^{(n+1)^c} = m$.
\end{lemma}

\begin{lemma}[Failure excludes every split]
\lean{Complexity.a3_split_none_iff}
\label{Complexity.a3_split_none_iff}
\leanok
\uses{Complexity.a3Split}
\difficulty{3}
The exponential-split search on $(C, c, m)$ fails if and only if no natural
number $n$ satisfies $n + C \cdot 2^{(n+1)^c} = m$.
\end{lemma}

\begin{lemma}[Completeness of the split search]
\lean{Complexity.a3_split_complete}
\label{Complexity.a3_split_complete}
\leanok
\uses{Complexity.a3Split, Complexity.a3_split_none_iff, Complexity.a3_split_spec, Complexity.a3_split_strictMono}
\difficulty{3}
If $n + C \cdot 2^{(n+1)^c} = m$, then the exponential-split search on
$(C, c, m)$ returns exactly $n$.
\end{lemma}

\begin{lemma}[The split search rejects length zero]
\lean{Complexity.a3_split_empty}
\label{Complexity.a3_split_empty}
\leanok
\uses{Complexity.a3Split, Complexity.a3_split_none_iff}
\difficulty{2}
For every positive coefficient $C$ and every degree $c$, the exponential-split
search on $(C, c, 0)$ returns the failure value, so length zero admits no
exponential split.
\end{lemma}

\begin{definition}[The emitted split word]
\lean{Complexity.a3SplitWord}
\label{Complexity.a3SplitWord}
\leanok
\uses{Complexity.a3Split, Turing.pairEncode}
Given constants $C, c$ and a word $y$, the split word of $y$ is the paired encoding
of the first $i$ symbols of $y$ and the remaining suffix when the exponential-split
search on $(C, c, |y|)$ returns $i$, and is the empty word, signalling failure,
when the search fails.
\end{definition}

\begin{lemma}[Linear length of the split word]
\lean{Complexity.a3_split_length}
\label{Complexity.a3_split_length}
\leanok
\uses{Complexity.a3Split, Complexity.a3SplitWord, Complexity.a3_split_spec, Turing.pairEncode}
\difficulty{3}
For every word $y$, the emitted split word for $y$ has length at most
$2|y| + 2$, including on malformed (unsplittable) inputs.
\end{lemma}

\begin{lemma}[Polynomial-time computation of the split word]
\lean{Complexity.a3_split_timed}
\label{Complexity.a3_split_timed}
\leanok
\uses{Complexity.a3Split, Complexity.a3SplitWord, Complexity.a3_exp_bits_timed, Turing.FinTM, Turing.FinTM.ComputesFunInTime, Turing.FinTM.ComputesInTime, Turing.FinTM.ComputesInTime.mono, Turing.FinTM.computesFunInTime_splitSolveWith, Turing.solveSplitWith}
\difficulty{5}
For all constants $C, c$ there exist a finite machine and a constant $A$
computing the split-word function (the paired prefix/suffix encoding, or the
empty failure word) within time $A (n + 1)^{c+2}$ on inputs of length $n$.
\end{lemma}

\begin{definition}[The padding verifier]
\lean{Complexity.a3Verifier}
\label{Complexity.a3Verifier}
\leanok
\uses{Complexity.a3Split}
For a language $L$ over the binary alphabet and a degree $c$: the language of
words $y$ such that the exponential-split search with coefficient one succeeds
on $|y|$ with some prefix length $n$ and the first $n$ symbols of $y$ form a
word in $L$. Words whose length admits no split are rejected; certificate bits
carry no information.
\end{definition}

\begin{lemma}[The padding verifier on exact-width concatenations]
\lean{Complexity.a3_verifier_append}
\label{Complexity.a3_verifier_append}
\leanok
\uses{Complexity.a3Split, Complexity.a3Verifier, Complexity.a3_split_complete}
\difficulty{2}
If $|u| = 2^{(|x|+1)^c}$, then the concatenation $x \mathbin{+\!\!+} u$ belongs
to the padding verifier of $L$ at degree $c$ if and only if $x \in L$.
\end{lemma}

\begin{lemma}[The split equation bounds the source budget]
\lean{Complexity.a3_source_budget}
\label{Complexity.a3_source_budget}
\leanok
\uses{Complexity.a3Split, Complexity.a3_split_spec}
\difficulty{3}
If the exponential-split search with coefficient one returns $n$ on a length
$m$, then $a \cdot 2^{n^c} \le a \cdot m$ for every constant $a$, including
degree zero.
\end{lemma}

\begin{lemma}[The padding verifier is polynomial-time]
\lean{Complexity.a3_verifier_mem_P}
\label{Complexity.a3_verifier_mem_P}
\leanok
\uses{Complexity.P, Complexity.a3NonemptyTM, Complexity.a3Split, Complexity.a3SplitWord, Complexity.a3Verifier, Complexity.a3_comp_at, Complexity.a3_decider_clean, Complexity.a3_nonempty_computes, Complexity.a3_nonempty_nil, Complexity.a3_source_budget, Complexity.a3_split_length, Complexity.a3_split_spec, Complexity.a3_split_timed, Complexity.mem_P_iff, Turing.FinTM, Turing.FinTM.ComputesInTime, Turing.FinTM.ComputesInTime.mono, Turing.FinTM.DecidesInTime, Turing.FinTM.bufferedCompTM, Turing.FinTM.computesFunInTime_pairFst, Turing.MultiTapeTM.indicator, Turing.pairDecode_pairEncode, Turing.pairEncode}
\difficulty{6}
If some finite machine decides a language $L$ over the binary alphabet within
time $a \cdot 2^{n^c}$, then the padding verifier of $L$ at degree $c$ belongs
to $\mathrm{P}$.
\end{lemma}

\begin{theorem}[$\mathrm{EXP} \subseteq \mathrm{NEXP}$]
\lean{Complexity.EXP_subset_NEXP}
\label{Complexity.EXP_subset_NEXP}
\leanok
\uses{Complexity.EXP, Complexity.NEXP, Complexity.a3Verifier, Complexity.a3_verifier_append, Complexity.a3_verifier_mem_P}
\difficulty{3}
Every language in $\mathrm{EXP}$ lies in $\mathrm{NEXP}$: a language over the
binary alphabet decidable in time $2^{n^c}$ admits certificates of length
exactly $2^{(n+1)^c}$ together with a polynomial-time verifier language.
\end{theorem}
